% \iffalse meta-comment
% !TeX program  = XeLaTeX
% !TeX encoding = UTF-8
%
% Copyright (C) 2021--2026 by Koyamin <contact@koyamin.com>
%
% This work may be distributed and/or modified under the
% conditions of the LaTeX Project Public License, either
% version 1.3c of this license or (at your option) any later
% version. The latest version of this license is in:
%
%   http://www.latex-project.org/lppl.txt
%
% and version 1.3 or later is part of all distributions of
% LaTeX version 2005/12/01 or later.
% 
% \fi
%
% \iffalse
%<class|ecnudoc|ecnulogo>\NeedsTeXFormat{LaTeX2e}
%<class|ecnudoc|ecnulogo>\RequirePackage{expl3}
%<class|ecnudoc|ecnulogo|definition>\GetIdInfo $Id: ecnuthesis.dtx 0.4 2026-04-20 15:00:00Z Koyamin <contact@koyamin.com>$
%<class>  {LaTeX Thesis Template for East China Normal University (ECNU)}
%<ecnudoc>  {Documentation class for ecnuthesis}
%<ecnulogo>  {logo for ecnuthesis}
%<class>\ProvidesExplClass{\ExplFileName}
%<ecnudoc>\ProvidesExplClass{ecnudoc}
%<ecnulogo>\ProvidesExplPackage{ecnulogo}
%<definition>  {Definition file for ecnuthesis}
%<definition>\ProvidesExplFile{\ExplFileName.def}
%<class|ecnudoc|ecnulogo|definition>  {\ExplFileDate}{\ExplFileVersion}{\ExplFileDescription}
%<*driver>
\documentclass{ecnudoc}
\EnableCrossrefs
\CodelineIndex
\RecordChanges
\begin{document}
  \DocInput{ecnuthesis.dtx}
  \IndexLayout
  \PrintIndex
\end{document}
%</driver>
% \fi
%
% \title{\textcolor{color-ecnu-red}{\textbf{ecnuthesis: 华东师范大学学位论文 \LaTeX 模板}}}
% \author{Koyamin}
% \date{v0.4\quad(2026-04-20)}
% \maketitle
% \tableofcontents
% \clearpage
%
% \begin{documentation}
%
% \section{模板介绍}
%
% \cls{ecnuthesis} （\textbf{E}ast \textbf{C}hina 
% \textbf{N}ormal \textbf{U}niversity \LaTeX{} \textbf{Thesis} 
% Template）是华东师范大学学位论文模板的 \LaTeX 实现，
% 旨在帮助华东师范大学的学生使用 \LaTeX{}
% 完成博士、硕士学位论文与本科生毕业论文的排版工作。
%
% 本模板根据华东师范大学研究生院发布的
% 《\href{https://yjsy.ecnu.edu.cn/6e/bf/c42090a683711/page.htm}{华东师范大学博士、硕士学位论文基本格式要求}》
% （发布于 2025 年 4 月 25 日）与
% 《\href{https://bksy.ecnu.edu.cn/d4/be/c40573a513214/page.htm}{华东师范大学本科生毕业论文（设计）格式要求}》
% （发布于 2021 年 11 月 10 日）制作而成，实现了包括
% 内封面、声明页、目录、摘要、正文、参考文献、附录、致谢等要素。
% 
%
% \section{声明}
%
% 本模版是华东师范大学学位论文模板的一个第三方 \LaTeX{} 开源实现。
% 本模板的发布遵守 \LaTeX{} Project Public License (1.3c) 协议。
% 使用本模板前，请务必阅读并同意以下事项：
%
% \begin{enumerate}
%   \item 本模版的开发行为未得到华东师范大学相关部门授权，本模版未经华东师范大学
%     相关部门审核。任何因使用本模版而导致的论文格式审查问题与本模版作者无关。
%   \item 使用本模板即视为已阅读并同意 \LaTeX{} Project Public License (1.3c) 
%     协议的全部内容。
%   \item 任何个人或组织在修改、扩展本模板并生成新的专用模板时，需严格遵守
%     \LaTeX{} Project Public License (1.3c) 协议。
%     因违反协议而导致的任何纠纷争端与本模板作者无关。
% \end{enumerate}
%
%
% \section{模板的获取与使用}
%
% \subsection{获取模板}
%
% 本模板的源代码托管在 GitHub 仓库 
% \href{https://github.com/koyamin/ecnuthesis}{Koyamin/ecnuthesis} 中。
% 用户可通过以下方式获得本模板：
%
% \begin{itemize}
%   \item GitHub 仓库的 Release 中提供了开箱即用的模板，仅需从 Release 中下载 zip 文件即可。
%   将下载至本地的 zip 文件解压缩后，打开解压缩得到的文件夹，
%   使用该文件夹目录中的 \file{thesis.tex} 来编写论文。
%
%   \item 本模板提供了 \href{https://www.overleaf.com/latex/templates/ecnuthesis-latex-thesis-template-for-east-china-normal-university/szppdtkvgvpk}{Overleaf 版本}，
%   打开链接并登录后即可直接编辑。
%
%   \item 可通过源代码 \file{ecnuthesis.dtx} 与 \file{ecnuthesis.ins} 生成
%   模板文档类 \file{ecnuthesis.cls} 文件。需在项目目录下运行以下命令生成
%   \file{ecnuthesis.cls} 文件：
% \begin{shellcode}
%   $ xetex ecnuthesis.ins
% \end{shellcode}
%   待 \file{ecnuthesis.cls} 文件生成完毕后，即可载入该文档类撰写论文。
% \end{itemize}
%
%
% \subsection{模版组成}
%
% \cls{ecnuthesis} 的主要组成部分由表
% \ref{table:ecnuthesis-components} 所示。
%
% \begin{table}[htb]
%   \caption{\cls{ecnuthesis} 的主要组成部分}
%   \label{table:ecnuthesis-components}
%   \centering \small
% \begin{tabular}{lll}
%   \toprule
%   \textbf{类别} & \textbf{文件} & \textbf{说明} \\ \midrule
%   开发文件 & \file{ecnuthesis.dtx} & DOCSTRIP 源文件  \\
%           & \file{ecnuthesis.ins} & DOCSTRIP 驱动文件   \\ \midrule
%   文档类 & \file{ecnuthesis.cls} & \cls{ecnuthesis} 文档类文件 \\
%         & \file{ecnuthesis.def} & 文档类名称配置 \\ \midrule
%   视觉形象系统 & \file{ecnu-vi-inner-cover-logo-ug.pdf} & 本科生毕业论文内封面校徽 \\
%               & \file{ecnu-vi-inner-cover-logo-pg.pdf} & 博士、硕士学位论文内封面校徽 \\ \midrule
%   用户手册 & \file{ecnuthesis.pdf} & 模板用户手册（本文档） \\
%   \bottomrule
% \end{tabular}
% \end{table}
% 
% 其中，文档类与视觉形象系统是本模板中不可缺少的两部分。
% 若缺少其中之一，则可能无法生成一份格式正确的论文。
%
% \subsection{使用模板编写论文}
%
% 以下内容需要用户具有处理一般的 \LaTeX 文档的能力。如果从未接触
% 过 \TeX{} 和 \LaTeX{}，建议先学习相关的基础知识。
%
%
% 以下是一份 \TeX 源文件，该文件展示了 \cls{ecnuthesis} 文档类最基本的用法。
%
% \begin{latexcode}[moretexcs={\ecnusetup,\tableofcontents,\chapter},
%   moretexcs={[2]\frontmatter,\mainmatter},
%   emph={[1]ecnuthesis,document},
%   emph={[2]type,oneside,style,info},
%   emph={[3]font,font-cjk,title,title-en,author,department,major}]
%   % thesis.tex
%   \documentclass[type=bachelor,oneside]{ecnuthesis}
%   \ecnusetup{
%     info/title = {计算机器与智能},
%     info/author = {艾伦·麦席森·图灵},
%     style/font-cjk = fandol,
%   }
%   \begin{document}
%     \frontmatter
%     \tableofcontents
%     \mainmatter
%     \chapter{欢迎}
%     \section{欢迎使用 ECNUThesis!}
%     你好, \LaTeX{}!
%   \end{document}
% \end{latexcode}
%
%
% \subsection{编译文档}\label{subsection:compile-doc}
%
% 本模板仅支持使用 \XeLaTeX 引擎编译。为了生成正确的目录、脚注以及交叉引用，
% 您至少需要连续编译两次。
%
% 假设编写论文所使用的 \TeX 源文件名为 \file{thesis.tex}。若要编译论文，请
% 在命令行中执行
% \begin{shellcode}
%   $ xelatex thesis
%   $ biber thesis    # 若引用了参考文献, 请执行该命令
%   $ xelatex thesis
% \end{shellcode}
% 或者使用 \pkg{latexmk}：
% \begin{shellcode}
%   $ latexmk -xelatex thesis
% \end{shellcode}
%
% \section{配置文档类}
%
% 为了使用户更方便地使用本模板编写论文，\cls{ecnuthesis} 提供
% 了若干易于使用的命令、环境和配置项，包括载入文档类时可配置的文档类选项、
% 配置论文信息与格式的导言区命令，以及编写论文时可能会使用到的命令与环境。
%
% 本模板所提供的命令、配置项或环境可以分为以下三类：
% \begin{itemize}
% \item 名字后带有 \rexptarget\rexpstar{} 标记的，表示只有在博士、硕士学位论文中方为有效；
% \item 名字后带有 \exptarget\expstar{} 标记的，表示只有在本科生毕业论文中方为有效；
% \item 名字后不带有特殊标记的，表示该选项在所有类型的论文中均为有效。
% \end{itemize}
% 后续文档将在使用说明中对某些特殊的选项加以说明。
%
% \subsection{文档类选项}\label{subsection:template-option}
%
% 文档类选项是指在载入文档类的时候需要指定的选项。载入文档类时，
% 用户可通过指定文档类选项来设置论文最基本的要素，包括论文类型、单双面
% 模式、是否生成封面以及声明页等。
% 
% \begin{latexcode}[emph={[1]ecnuthesis}]
%   \documentclass(*\oarg{文档类选项}*){ecnuthesis}
% \end{latexcode}
% 
% 文档类选项的形式通常为 \kvopt{\meta{key}}{\meta{value}}，
% 不同的文档类选项之间使用半角逗号（即“|,|”）分隔。
% 一些文档类选项为布尔型，它们只能在 |true| 和 |false|
% 中取值。对于这些选项，\kvopt{\meta{选项}}{true} 中的 |=true|
% 可以被省略。
%
% \begin{function}{type}
%   \begin{syntax}
%     \OPTKEY{type} = bachelor|master|doctor
%   \end{syntax}
%   选择论文学位类型。该选项有三个可选值，这三个可选值的含义分别为：
%   \begin{optdesc}
%     \item[bachelor] 表示本科生毕业论文。
%     \item[master] 表示硕士学位论文。
%     \item[doctor] 表示博士学位论文。
%   \end{optdesc}
% \end{function}
% 
% \begin{function}{oneside,twoside}
%   指定论文的排版模式。该选项会影响每章的开始位置，
%   还会影响页眉样式。模板默认的排版模式为双面模式（\opt{twoside}）。
%
%   在双面模式（\opt{twoside}）下，按照通常的排版惯例，每章应只从
%   奇数页（在右）开始；而在单页模式（\opt{oneside}）下，则可以从
%   任意页面开始。本模板中，目录、摘要等均视作章，也按相同方式排版。
%   使用双面模式适合双面打印纸质版论文；使用单面模式适合提交电子版论文。
% \end{function}
%
% \begin{function}{cover}
%   \begin{syntax}
%     \OPTKEY{cover} = (true)|false
%   \end{syntax}
%   选择是否生成封面页。模板默认自动生成封面页。
%   若设置 \opt{cover=true}，模版将自动生成封面页；反之则不生成封面页。
% \end{function}
%
% \begin{function}{declaration}
%   \begin{syntax}
%     \OPTKEY{declaration} = true|(false)
%   \end{syntax}
%   选择是否生成声明页。模板默认不自动生成声明页。
%   若缺省则不生成声明页，反之则生成声明页。
%   声明页默认由模版自动生成，也可以使用
%   \opt{style/declaration-src} 选项，
%   将指定的 PDF 文件作为声明页插入到文档中。
% 
%   若选择不生成声明页，则 \opt{style/declaration-src} 选项将失效。
% \end{function}
%
%
% \subsection{配置命令}\label{subsection:template-config-command}
%
% \cls{ecnuthesis} 提供了一个用于导言区的配置命令 \cs{ecnusetup}，用于
% 录入论文相关信息，并提供细粒度的格式设置。
%
% \begin{function}{\ecnusetup}
%   \begin{syntax}
%     \TNi{ecnusetup}\marg{键值对列表}
%   \end{syntax}
%   \cs{ecnusetup} 命令用于配置相关自定义选项。
%   载入 \cls{ecnuthesis} 文档类后，本小节中的所有选项均可使用
%   \cs{ecnusetup} 命令配置。
%   该命令须写在导言区中。
% \end{function}
%
% \cs{ecnusetup} 的参数是一组由半角逗号（即“|,|”）分隔的选项列表，
% 列表中的选项通常为 \kvopt{\meta{key}}{\meta{value}} 的形式，我们称这种
% 形式为键值，由键值组成的选项列表称为键值列表。键值中的 |=| 周围的
% 空格不影响设置，但是键值列表之间不可以出现多余的空行。对于同一个选项，后面的
% 设置将会覆盖前一个设置。
%
% 一些选项包含若干子选项（例如 \opt{style} 选项和 \opt{info} 选项）。
% 此时，可选择这两种等价的设置方式来配置论文：
% \begin{latexcode}[moretexcs={\ecnusetup},
%   emph={[2]type,oneside,style,info},
%   emph={[3]font,font-cjk,title,title-en,author,department,major}]
%   \ecnusetup{
%     style = { font = times, font-cjk = fandol },
%     info  = {
%       title       = {计算机器与智能},
%       title-en    = {Computing Machinery and Intelligence},
%       author      = {艾伦·麦席森·图灵},
%       department  = {计算机科学与技术学院},
%       major       = {计算机科学与技术}
%     }
%   }
% \end{latexcode}
% 或者
% \begin{latexcode}[moretexcs={\ecnusetup},
%   emph={[2]type,oneside,style,info},
%   emph={[3]font,font-cjk,title,title-en,author,department,major}]
%   \ecnusetup{
%     style/font      = times,
%     style/font-cjk  = fandol,
%     info/title      = {计算机器与智能},
%     info/title-en   = {Computing Machinery and Intelligence},
%     info/author     = {艾伦·麦席森·图灵},
%     info/department = {计算机科学与技术学院},
%     info/major      = {计算机科学与技术}
%   }
% \end{latexcode}
% 请注意，斜杠符号 |/| 的前后均不可以出现空白字符。
%
% \subsubsection{论文基本信息}
%
% 在 \cls{ecnuthesis} 中，用户可以通过配置 \opt{info} 选项
% 的子选项来完成论文基本信息的填写。本章节将介绍 \opt{info} 选项的所有可供
% 配置的子选项。
%
% \begin{function}{info}
%   \begin{syntax}
%     \CFG{info} = \marg{键值对列表}
%     \CFG{info}/\meta{key} = \meta{value}
%   \end{syntax}
%   该选项用于录入论文信息。该选项拥有一些子选项，具体的设置子选项如下所示。
% \end{function}
%
% \begin{function}[rEXP]{info/degree}
%   \begin{syntax}
%     \CFGKEY{degree} = (academic)|professional
%   \end{syntax}
%   博士、硕士学位论文的学位类型。
%   \begin{optdesc}
%     \item[academic] 表示学术学位。
%     \item[professional] 表示专业学位。
%   \end{optdesc}
% \end{function}
%
% \begin{function}[rEXP]{info/clc}
%   \begin{syntax}
%     \CFGKEY{clc} = \marg{中图分类号}
%   \end{syntax}
%   中图分类号（CLC）。
% \end{function}
%
% \begin{function}[rEXP]{info/secret-level}
%   \begin{syntax}
%     \CFGKEY{secret-level} = \marg{密级}
%   \end{syntax}
%   密级。
% \end{function}
%
% \begin{function}{info/title,info/title-en}
%   \begin{syntax}
%     \CFGKEY{title} = \marg{中文标题}
%     \CFGKEY{title-en} = \marg{外文标题}
%   \end{syntax}
%   论文的中外文标题。可以在标题内部使用换行命令 |\\|。
% \end{function}
%
% \begin{function}{info/author}
%   \begin{syntax}
%     \CFGKEY{author} = \marg{中文姓名}
%   \end{syntax}
%   作者的中文姓名。
% \end{function}
%
% \begin{function}[rEXP]{info/author-en}
%   \begin{syntax}
%     \CFGKEY{author-en} = \marg{外文姓名}
%   \end{syntax}
%   作者的外文姓名。
% \end{function}
% 
% \begin{function}{info/student-id}
%   \begin{syntax}
%     \CFGKEY{student-id} = \marg{学号}
%   \end{syntax}
%   作者的学号。
% \end{function}
%
% \begin{function}{info/department}
%   \begin{syntax}
%     \CFGKEY{department} = \marg{院系中文名称}
%   \end{syntax}
%   作者所在院系的中文名称。
% \end{function}
%
% \begin{function}[rEXP]{info/department-en}
%   \begin{syntax}
%     \CFGKEY{department-en} = \marg{院系外文名称}
%   \end{syntax}
%   作者所在院系的外文名称。
% \end{function}
%
% \begin{function}{info/major}
%   \begin{syntax}
%     \CFGKEY{major} = \marg{专业中文名称}
%   \end{syntax}
%   作者所在专业的中文名称。
% \end{function}
%
% \begin{function}[rEXP]{info/major-en}
%   \begin{syntax}
%     \CFGKEY{major-en} = \marg{专业外文名称}
%   \end{syntax}
%   作者所在专业的外文名称。
% \end{function}
%
% \begin{function}[rEXP]{info/field,info/field-en}
%   \begin{syntax}
%     \CFGKEY{field} = \marg{研究方向中文名称}
%     \CFGKEY{field-en} = \marg{研究方向外文名称}
%   \end{syntax}
%   作者所在研究方向的中外文名称。用于填写博士、硕士毕业论文内封面
%   中的“研究方向”信息（学术学位）或“专业学位领域”信息（专业学位）。
% \end{function}
%
% \begin{function}{info/supervisor}
%   \begin{syntax}
%     \CFGKEY{supervisor} = \marg{指导教师中文姓名}
%   \end{syntax}
%   论文指导教师姓名。用于填写内封面的“指导教师”信息。
% \end{function}
%
% \begin{function}[rEXP]{info/supervisor-en}
%   \begin{syntax}
%     \CFGKEY{supervisor-en} = \marg{指导教师外文姓名}
%   \end{syntax}
%   论文指导教师姓名。用于填写内封面的“指导教师”信息。
% \end{function}
%
% \begin{function}[EXP]{info/academic-title}
%   \begin{syntax}
%     \CFGKEY{academic-title} = \marg{职称}
%   \end{syntax}
%   论文指导教师职称。用于填写内封面的“职称”信息。
% \end{function}
%
% \begin{function}{info/year}
%   \begin{syntax}
%     \CFGKEY{year} = \meta{论文完成年份}
%   \end{syntax}
%   论文完成年份。用于填写内封面底部的年份信息。
%   若该选项被省略，模版将调用当前年份。请注意，该选项没有花括号。
% \end{function}
%
% \begin{function}{info/month}
%   \begin{syntax}
%     \CFGKEY{month} = \meta{论文完成月份}
%   \end{syntax}
%   论文完成月份。用于填写内封面底部的月份信息。
%   若该选项被省略，模版将调用当前月份。请注意，该选项没有花括号。
% \end{function}
%
% \begin{function}[EXP]{info/graduation-year}
%   \begin{syntax}
%     \CFGKEY{graduation-year} = \meta{毕业年份}
%   \end{syntax}
%   内封面毕业届别。用于填写内封面左上角的届别信息。
%   若该选项被省略，模版将调用 \opt{year} 字段的值；
%   若 \opt{year} 选项也被省略，模版将调用当前年份。
%   请注意，该选项没有花括号。
% \end{function}
%
% \begin{function}{info/keywords,info/keywords-en}
%   \begin{syntax}
%     \CFGKEY{keywords} = \marg{中文关键词}
%     \CFGKEY{keywords-en} = \marg{外文关键词}
%   \end{syntax}
%   中外文关键词列表，用于填写中外文摘要页的关键词信息。
%   请使用半角逗号（即“|,|”）分隔关键词。
% \end{function}
%
% \subsubsection{论文格式设置}
%
% 在 \cls{ecnuthesis} 中，用户可以通过配置 \opt{style} 选项
% 的子选项来完成论文格式方面的相关设置。本章节将介绍 \opt{style} 选项的所有可供
% 配置的子选项。
%
% \begin{function}{style}
%   \begin{syntax}
%     \CFG{style} = \marg{键值对列表}
%     \CFG{style}/\meta{key} = \meta{value}
%   \end{syntax}
%   该选项用于设置论文格式相关选项。其拥有一些子选项，具体的设置子选项如下所示。
% \end{function}
%
% \begin{function}{style/numbering}
%   \begin{syntax}
%     \CFGKEY{numbering} = (arabic)|chinese
%   \end{syntax}
%   设置论文的章节编号样式。
%   \begin{optdesc}
%     \item[arabic] 使用阿拉伯数字进行编号（即理科要求）。
%     \item[chinese] 使用中文数字进行编号（即文科要求）。
%   \end{optdesc}
% \end{function}
%
% \begin{function}{style/font-cjk}
%   \begin{syntax}
%     \CFGKEY{font-cjk} = (auto)|fandol|mac|windows
%   \end{syntax}
%   设置论文的中文字体。具体的各配置选项及其所对应的中文字体
%   见表 \ref{table:font-cjk}。默认选项为 \opt{auto}，
%   此时会根据操作系统自动配置自动配置：Windows 系统默认使用 \opt{windows}，
%   macOS 系统默认使用 \opt{mac}，Linux 系统默认使用 \opt{fandol}。
% \end{function}
%
% \begin{table}[ht]
%   \caption{各中文字体配置选项所对应的中文字体}
%   \label{table:font-cjk}
%   \centering \small
%   \begin{tabular}{ccccc}
%   \toprule
%           & \textbf{正文字体（宋体）} & \textbf{无衬线字体（黑体）} 
%           & \textbf{等宽字体（仿宋）} & \textbf{楷体}       \\
%   \midrule
%   \opt{windows} & 中易宋体     & 中易黑体      & 中易仿宋     & 中易楷体     \\
%   \opt{mac}     & 华文宋体-简   & 华文黑体-简    & 华文仿宋     & 华文楷体-简   \\
%   \opt{fandol}  & Fandol宋体 & Fandol黑体  & Fandol仿宋 & Fandol楷体  \\
%   \bottomrule
%   \end{tabular}
% \end{table}
%
% \begin{function}{style/font}
%   \begin{syntax}
%     \CFGKEY{font} = times|(xits)|lm
%   \end{syntax}
%   设置论文的西文和数学字体。默认选项为 \opt{xits}。具体的各配置选项及其所对应的西文字体
%   与数学字体见表 \ref{table:font} 。
% \end{function}
%
% \begin{table}[ht]
%   \caption{各西文字体配置选项所对应的西文字体}
%   \label{table:font}
%   \centering \small
%   \begin{tabular}{ccccc}
%   \toprule
%       & \textbf{衬线字体} & \textbf{无衬线字体} & \textbf{等宽字体} & \textbf{数学字体} \\
%   \midrule
%   \opt{times} & Times New Roman & Arial    & Courier New & XITS Math \\
%   \opt{xits}  & XITS            & TG Heros & TG Cursor   & XITS Math \\
%   \opt{lm}    & LM Roman        & LM Sans  & LM Mono     & LM Math  \\
%   \bottomrule
%   \end{tabular}
% \end{table}
%
% \begin{function}{style/bib-src}
%   \begin{syntax}
%     \CFGKEY{bib-src} = \marg{文件}
%   \end{syntax}
%   设置参考文献数据源文件。此处需填写数据源文件所在路径（可以是绝对路径或相对路径）。
%   由于参考文献使用 biber 进行处理，故必须明确给出该数据源文件的
%   \file{.bib} 后缀名。
% \end{function}
%
% \begin{function}{style/logo-src}
%   \begin{syntax}
%     \CFGKEY{logo-src} = \marg{文件}
%   \end{syntax}
%   设置封面插图数据源。此处需填写数据源文件所在路径（可以是绝对路径或相对路径）。
%   模版已自带封面插图数据源，位于 \file{./source/inner-cover(contains\_font).eps}。
% \end{function}
%
% \begin{function}{style/declaration-src}
%   \begin{syntax}
%     \CFGKEY{declaration-src} = \marg{PDF 文件}
%   \end{syntax}
%   设置待插入的扫描版声明页 PDF 文件。若该值为空，则生成模版预定义的声明页；
%   否则将指定路径所对应的 PDF 文件作为声明页插入。
% \end{function}
%
%
% \section{编写论文内容}
%
% 为符合论文的格式要求，\cls{ecnuthesis} 文档类提供了一系列命令与环境，
% 方便用户完成学位论文内容的撰写工作。
%
% 一篇学位论文的内容通常由以下三个部分组成：
% \begin{itemize}
%   \item 前置部分：中外文摘要、目录（包括插图目录和表格目录）；
%   \item 正文部分：各个章节；
%   \item 后置部分：参考文献、附录、后记（致谢）。
% \end{itemize}
% 本节将根据以上三个部分，分别介绍 \cls{ecnuthesis} 文档类所提供的
% 相关命令与环境。
%
% \subsection{前置部分}\label{subsubsection:frontmatter}
%
% 前置部分包括中外文摘要、目录（包括插图目录和表格目录）。
%
% \begin{function}{\frontmatter}
%   \begin{syntax}
%     \TNii{frontmatter}
%   \end{syntax}
%   用于声明论文前置部分开始。
%   使用 \cs{frontmatter} 命令后，前置部分的页码采用大写罗马字母，
%   并且与正文分开计数。
% \end{function}
%
% \begin{function}{\tableofcontents,\listoftables,\listoffigures}
%   用于生成目录、插图目录以及表格目录。
% \end{function}
%
% \begin{function}{abstract,abstractEN}
%   \begin{syntax}
%     \BEV{abstract}
%     \  \meta{中文摘要}
%     \EEV{abstract}
%     \BEV{abstractEN}
%     \  \meta{外文摘要}
%     \EEV{abstractEN}
%   \end{syntax}
%   中外文摘要页环境。该环境会在摘要页的末尾自动添加关键词。
% \end{function}
%
% \subsection{正文部分}\label{subsubsection:mainmatter}
%
% \begin{function}{\mainmatter}
%   \begin{syntax}
%     \TNii{mainmatter}
%   \end{syntax}
%   用于声明论文正文部分开始。使用 \cs{mainmatter} 命令后，
%   正文部分的页码采用阿拉伯数字。
% \end{function}
%
% \begin{function}[EXP]{\bicaption}
%   \begin{syntax}
%     \TNi{bicaption}\marg{中文题注}\marg{外文题注}
%   \end{syntax}
%   本科生毕业论文中插图和表格的中外文对照题注。
%   与 \cs{caption} 命令类似，仅需在 \env{figure} 和 \env{table}
%   环境中的合适位置编写该命令即可。
%   以下是一个使用中外文对照题注的示例：
% \begin{latexcode}[moretexcs={\includegraphics,\bicaption},emph={[1]figure}]
%   \begin{figure}[htb]
%     \centering
%     \includegraphics[width=.5\textwidth]{component-deployment.png}
%     \bicaption{组件分布图}{Component Deployment}
%   \end{figure}
% \end{latexcode}
% \end{function}
%
%
% \subsection{后置部分}\label{subsubsection:backmatter}
%
% 后置部分包括参考文献列表、附录以及致谢等。
%
% \begin{function}{\backmatter}
%   \begin{syntax}
%     \TNii{backmatter}
%   \end{syntax}
%   用于声明论文后置部分开始。
% \end{function}
%
% \begin{function}{\printbibliography}
%   \begin{syntax}
%     \TNi{printbibliography}
%   \end{syntax}
%   用于生成参考文献列表。
% \end{function}
%
% \begin{function}{appendix}
%   \begin{syntax}
%     \BEV{appendix}
%     \  \meta{附录}
%     \EEV{appendix}
%   \end{syntax}
%   附录环境。
% \end{function}
%
% \begin{function}{acknowledgement}
%   \begin{syntax}
%     \BEV{acknowledgement}
%     \  \meta{致谢}
%     \EEV{acknowledgement}
%   \end{syntax}
%   致谢环境。
% \end{function}
%
% \section{宏包依赖情况}
%
% \cls{ecnuthesis} 在任何情况下均会显式调用以下宏包或文档类：
% \begin{itemize}
%   \item \pkg{xparse} 宏包、\pkg{xtemplate} 宏包以及 \pkg{l3keys2e} 宏包，
%     用于扩展 \LaTeX3 编程环境。属于 \pkg{l3packages} 宏包。
%   \item \cls{ctexrep} 文档类，提供中文排版的通用框架。属于 \CTeX{} 宏包。
%   \item \pkg{amsmath} 宏包，用于增强 \LaTeX{} 的数学排版功能。
%   \item \pkg{unicode-math} 宏包，用于处理使用 Unicode 编码的 OpenType 数学字体。
%   \item \pkg{geometry} 宏包，用于调整页面尺寸。
%   \item \pkg{fancyhdr} 宏包，用于设置页眉和页脚。
%   \item \pkg{footmisc} 宏包，用于设置脚注。
%   \item \pkg{caption} 宏包，用于设置图片与表格的题注。
%   \item \pkg{bicaption} 宏包，用于支持中外文对照版题注。
%   \item \pkg{xcolor} 宏包，用于提供彩色支持。
%   \item \pkg{xeCJKfntef} 宏包，用于实现内封面页的中文下划线。
%   \item \pkg{titletoc} 宏包，用于设置目录格式。
%   \item \pkg{biblatex} 宏包，用于参考文献的格式设置，并依赖 biber 程序。
%   \item \pkg{hyperref} 宏包，用于提供交叉引用、超链接以及电子书签等功能。
% \end{itemize}
%
% 此外，若指定不同的选项，模板可能会调用其他的宏包。具体而言：
% \begin{itemize}
%   \item 若 \cs{ecnusetup} 命令的 \opt{style/declaration-src} 选项指定了
%   需要插入至论文的 PDF 声明页文件，则本模板还将调用 \pkg{pdfpages} 宏包。
% \end{itemize}
%
%
%
% \end{documentation}
%
% 
%
% \begin{implementation}
%
% 
%
% \section{实现细节}
%
% \cls{ecnuthesis} 使用 \LaTeX3 语法编写，
% 依赖 \pkg{expl3} 环境，并依赖 \pkg{l3packages} 中的相关宏包。
%
% 根据 \LaTeX3 的语法规则，代码中的空格、换行、回车与制表符将被完全忽略，
% 而下划线 |\_| 和冒号 |:| 则可作为一般字母使用。
%
% \subsection{准备}
%
% 
%    \begin{macrocode}
%<@@=ecnu>
%<*class>
%    \end{macrocode}
%
% 导入 \LaTeX3 编程环境。
%
%    \begin{macrocode}
\RequirePackage { xparse, xtemplate, l3keys2e }
%    \end{macrocode}
%
% 目前 \cls{ecnuthesis} 仅支持 \XeTeX{}。
%    \begin{macrocode}
\msg_new:nnn { ecnuthesis } { engine-unsupportd } 
  { \\
    This~ class~ requires~ XeTeX~ to~ compile. \\
    The~ engine~ "#1"~ is~ not~ supported. 
  }
\sys_if_engine_xetex:F
  {
    \msg_fatal:nnx { ecnuthesis } { engine-unsupportd } { \c_sys_engine_str }
  }
%    \end{macrocode}
%
% \subsubsection{内部变量声明}
%
% \begin{variable}{\l_@@_tmpa_box,
%   \l_@@_tmpb_box,
%   \l_@@_tmpc_box,
%   \l_@@_tmpa_clist,
%   \l_@@_tmpb_clist,
%   \l_@@_tmpa_skip,
%   \l_@@_tmpa_tl,
%   \l_@@_tmpb_tl,
%   \l_@@_tmpa_dim,
%   \l_@@_tmpb_dim,
%   \l_@@_tmpc_dim}
% 声明一些临时变量。
%    \begin{macrocode}
\box_new:N   \l_@@_tmpa_box
\box_new:N   \l_@@_tmpb_box
\box_new:N   \l_@@_tmpc_box
\clist_new:N \l_@@_tmpa_clist
\clist_new:N \l_@@_tmpb_clist
\skip_new:N  \l_@@_tmpa_skip
\tl_new:N    \l_@@_tmpa_tl
\tl_new:N    \l_@@_tmpb_tl
\dim_new:N   \l_@@_tmpa_dim
\dim_new:N   \l_@@_tmpb_dim
\dim_new:N   \l_@@_tmpc_dim
%    \end{macrocode}
% \end{variable}
%
% \begin{variable}{\g_@@_thesis_type_tl, \g_@@_is_bachelor_bool}
% 论文类型。用于判断论文是否为本科生毕业论文。
%    \begin{macrocode}
\tl_new:N \g_@@_thesis_type_tl
\bool_new:N \g_@@_is_bachelor_bool
\bool_set_false:N \g_@@_is_bachelor_bool
%    \end{macrocode}
% \end{variable}
%
% \begin{variable}{\g_@@_twoside_bool}
% 是否开启双面打印模式（默认关闭）。
%    \begin{macrocode}
\bool_new:N \g_@@_twoside_bool
\bool_set_false:N \g_@@_twoside_bool
%    \end{macrocode}
% \end{variable}
%
% \begin{variable}{\g_@@_cover_page_bool}
% 是否自动生成封面页（默认自动生成）。
%    \begin{macrocode}
\bool_new:N \g_@@_cover_page_bool
\bool_set_true:N \g_@@_cover_page_bool
%    \end{macrocode}
% \end{variable}
%
% \begin{variable}{\g_@@_declaration_page_bool}
% 是否自动生成声明页（默认不自动生成）。
%    \begin{macrocode}
\bool_new:N \g_@@_declaration_page_bool
\bool_set_false:N \g_@@_declaration_page_bool
%    \end{macrocode}
% \end{variable}
%
% \begin{variable}{\g_@@_to_ctexrep_clist}
% 保存由 \cls{ecnuthesis} 传入 \cls{ctexrep} 宏包
% 的选项列表。
%    \begin{macrocode}
\clist_new:N \g_@@_to_ctexrep_clist
%    \end{macrocode}
% \end{variable}
%
% \subsubsection{内部函数}
%
% \begin{macro}{\@@_quad:,\@@_qquad:}
% 等价于 \LaTeXe{} 中的 \tn{quad} 和 \tn{qquad}。
%    \begin{macrocode}
\cs_new:Npn \@@_quad:  { \skip_horizontal:n { 1 em } }
\cs_new:Npn \@@_qquad: { \skip_horizontal:n { 2 em } }
%    \end{macrocode}
% \end{macro}
%
% \begin{macro}{\@@_vspace:N,\@@_vspace:n,\@@_vspace:c}
% 类似 \LaTeXe{} 中的 \tn{vspace*}。
%    \begin{macrocode}
\cs_new_protected:Npn \@@_vspace:N #1
  {
    \dim_set_eq:NN \l_@@_tmpa_dim \prevdepth
    \hrule height \c_zero_dim
    \nobreak
    \skip_vertical:N #1
    \skip_vertical:N \c_zero_skip
    \dim_set_eq:NN \prevdepth \l_@@_tmpa_dim
  }
\cs_new_protected:Npn \@@_vspace:n #1
  {
    \skip_set:Nn \l_@@_tmpa_skip {#1}
    \@@_vspace:N \l_@@_tmpa_skip
  }
\cs_generate_variant:Nn \@@_vspace:N { c }
%    \end{macrocode}
% \end{macro}
%
% \begin{macro}{\@@_symbol:n}
% 等价于 \LaTeXe{} 中的 \tn{symbol}。
%    \begin{macrocode}
\cs_new:Npn \@@_symbol:n #1 { \tex_char:D #1 \scan_stop: }
%    \end{macrocode}
% \end{macro}
%
% \begin{macro}{\@@_arabic:n}
% 等价于 \LaTeXe{} 中的 \tn{arabic}。
%    \begin{macrocode}
\cs_new:Npn \@@_arabic:n #1
  { \int_to_arabic:v { c@ #1 } }
%    \end{macrocode}
% \end{macro}
%
% \begin{macro}{\@@_gadd_ltxhook:nn}
% 封装 \LaTeX{} 的钩子管理机制。本模板中的字体加载命令位于
% \opt{begindocument/\allowbreak before} 钩子中，需确保在 \pkg{xeCJK} 之前执行。
%    \begin{macrocode}
\cs_new_protected:Npn \@@_gadd_ltxhook:nn #1#2
  { \hook_gput_code:nnn {#1} { . } {#2} }
\hook_gset_rule:nnnn { begindocument/before } { . } { < } { xeCJK }
%    \end{macrocode}
% \end{macro}
%
% \begin{macro}{\@@_patch_cmd:Nnn,\@@_preto_cmd:Nnn,\@@_appto_cmd:Nn}
% 补丁工具，来自 \pkg{ctexpatch} 宏包。
%    \begin{macrocode}
\cs_new_protected:Npn \@@_patch_cmd:Nnn #1#2#3
  {
    \ctex_patch_cmd_once:NnnnTF #1 { } {#2} {#3}
      { } { \ctex_patch_failure:N #1 }
  }
\cs_new_protected:Npn \@@_preto_cmd:Nn #1#2
  {
    \ctex_preto_cmd:NnnTF #1 { } {#2}
      { } { \ctex_patch_failure:N #1 }
  }
\cs_new_protected:Npn \@@_appto_cmd:Nn #1#2
  {
    \ctex_appto_cmd:NnnTF #1 { } {#2}
      { } { \ctex_patch_failure:N #1 }
  }
%    \end{macrocode}
% \end{macro}
%
% \begin{macro}{\@@_msg_new:nn,
%   \@@_error:n,\@@_error:nn,\@@_error:nx,\@@_error:nnn,\@@_error:nnnn,
%   \@@_warning:n,\@@_warning:nn,\@@_warning:nxx,
%   \@@_info:nx}
% 各种信息函数的缩略形式。
%    \begin{macrocode}
\cs_new:Npn \@@_msg_new:nn  { \msg_new:nnn      { ecnuthesis } }
\cs_new:Npn \@@_error:n     { \msg_error:nn     { ecnuthesis } }
\cs_new:Npn \@@_error:nn    { \msg_error:nnn    { ecnuthesis } }
\cs_new:Npn \@@_error:nx    { \msg_error:nnx    { ecnuthesis } }
\cs_new:Npn \@@_error:nnn   { \msg_error:nnnn   { ecnuthesis } }
\cs_new:Npn \@@_error:nnnn  { \msg_error:nnnnn  { ecnuthesis } }
\cs_new:Npn \@@_warning:n   { \msg_warning:nn   { ecnuthesis } }
\cs_new:Npn \@@_warning:nn  { \msg_warning:nnn  { ecnuthesis } }
\cs_new:Npn \@@_warning:nxx { \msg_warning:nnxx { ecnuthesis } }
\cs_new:Npn \@@_info:nx     { \msg_info:nnx     { ecnuthesis } }
%    \end{macrocode}
% \end{macro}
%
% \subsection{选项处理}
%
% 定义 \opt{ecnu/option} 键值类。
%    \begin{macrocode}
\keys_define:nn { ecnu / option }
  {
%    \end{macrocode}
%
% \begin{macro}{type}
% 设置论文类型。默认为硕士论文。
%    \begin{macrocode}
    type        .choices:nn = 
      { doctor, master, bachelor }
      {
        \tl_set_eq:NN \g_@@_thesis_type_tl \l_keys_choice_tl
        \tl_if_eq:NnTF \l_keys_choice_tl { bachelor }
          { \bool_set_true:N \g_@@_is_bachelor_bool }
          { \bool_set_false:N \g_@@_is_bachelor_bool }
      },
    type .value_required:n  = true,
    type .initial:n         = master,
%    \end{macrocode}
% \end{macro}
%
% \begin{macro}{oneside,twoside}
% 设置开启打印模式。若开启打印模式，则将输出适合双面打印的文档，否则将输出适合单面打印的
% 文档。默认关闭打印模式。
%    \begin{macrocode}
    oneside .value_forbidden:n = true,
    twoside .value_forbidden:n = true,
    oneside .code:n =
      {
        \clist_gput_right:Nn \g_@@_to_ctexrep_clist { oneside }
        \bool_set_false:N    \g_@@_twoside_bool
      },
    twoside .code:n =
      {
        \clist_gput_right:Nn \g_@@_to_ctexrep_clist { twoside }
        \bool_set_true:N     \g_@@_twoside_bool
      },
%    \end{macrocode}
% \end{macro}
%
% \begin{macro}{cover}
% 设置开启封面页。若开启封面页，则模版将自动生成封面页；反之则不会自动生成封面页。
% 模板默认开启封面页。
%    \begin{macrocode}
    cover         .bool_gset:N = \g_@@_cover_page_bool,
    cover         .default:n = { true },
    cover         .initial:n = { true },
%    \end{macrocode}
% \end{macro}
%
% \begin{macro}{declaration}
% 设置开启声明页。若开启声明页，则模版将自动生成声明页；反之则不会自动生成声明页。
% 默认关闭声明页。
%    \begin{macrocode}
    declaration   .bool_gset:N = \g_@@_declaration_page_bool,
    declaration   .default:n = { true },
    declaration   .initial:n = { false },
%    \end{macrocode}
% \end{macro}
%
% 处理未知选项。
%    \begin{macrocode}
    unknown .code:n = { \@@_error:n { unknown-option } }
  }
\@@_msg_new:nn { unknown-option }
  { Class~ option~ "\l_keys_key_str"~ is~ unknown. }
%    \end{macrocode}
%
% 将文档类选项传给 \opt{ecnu/option}。
%    \begin{macrocode}
\ProcessKeysOptions { ecnu / option }
%    \end{macrocode}
%
% 载入预设文字
%    \begin{macrocode}
\file_input:n { ecnuthesis.def }
%    \end{macrocode}
%
% \subsection{行距}
%
% 由于导入 \cls{ctexrep} 文档类时需要指定行距，故在此先设置正文行距。
% 
% 根据《华东师范大学本科生毕业论文（设计）格式要求》的
% 相关要求，本科生毕业论文的正文字号应设置为小四（即 \qty{12}{bp}），
% 正文行间距应设置为 1.5 倍行距。
%
% 在 Word 中，默认的文档网格的行距是 \qty{15.6}{bp}。
% 当正文字号设置为小四时，若每一行对齐到该默认文档网格，则此时的行距
% 应为 $\qty{15.6}{bp} \times 1.5 = \qty{23.4}{bp}$。
% 在 \LaTeX{} 中，单倍行距的长度是字号长度的 1.2 倍。
% 这意味着 Word 中的 1.5 倍行距相当于 \LaTeX{} 中的
% $\qty{23.4}{bp} / \qty{12}{bp} / 1.2 = 1.625$ 倍行距。
% 故将本科生毕业论文的行距设置为
% \begin{macro}{\c_@@_ug_line_spread_fp}
% 本科生毕业论文的行距。
%    \begin{macrocode}
\fp_const:Nn \c_@@_ug_line_spread_fp { 1.625 }
%    \end{macrocode}
% \end{macro}
%
% 根据《华东师范大学博士、硕士学位论文基本格式要求》的
% 相关要求，博士、硕士学位论文的正文字号应设置为小四，
% 正文行间距应设置为 \qty{25}{bp}。所以，在 \LaTeX{} 中的行距应为
% $\qty{25}{bp} / (1.2\times\qty{12}{bp}) $。
% \begin{macro}{\c_@@_pg_line_spread_fp}
% 博士、硕士学位论文的行距。
%    \begin{macrocode}
\fp_const:Nn \c_@@_pg_line_spread_fp
  { \dim_ratio:nn { 25 bp } { 12 bp } / 1.2 }
%    \end{macrocode}
% \end{macro}
%
% 为了方便，此处再定义一个标准的单倍行距。
%    \begin{macrocode}
\fp_const:Nn \c_@@_single_space_line_spread_fp
  { \dim_ratio:nn { 15.6 bp } { 12 bp } / 1.2 }
%    \end{macrocode}
% 
% \subsection{载入文档类与宏包}
%
% 将选项传入 \cls{ctexrep} 文档类。
%    \begin{macrocode}
\PassOptionsToClass
  {
    UTF8,
    heading     = true,
    zihao       = -4,
    fontset     = none,
    linespread  = \bool_if:NTF \g_@@_is_bachelor_bool
      { \c_@@_ug_line_spread_fp }
      { \c_@@_pg_line_spread_fp },
    \g_@@_to_ctexrep_clist,
  }
  { ctexrep }
%    \end{macrocode}
%
% 将一些选项传入宏包。
%    \begin{macrocode}
\clist_map_inline:nn
  {
    { no-math           } { fontspec },
    { perpage           } { footmisc },
  }
  { \PassOptionsToPackage #1 }
%    \end{macrocode}
%
% 载入 \cls{ctexrep} 文档类。在使用 \XeTeX{} 编译时，\cls{ctexrep} 
% 的底层将调用 \pkg{xeCJK} 宏包以及 \pkg{fontspec} 宏包。
%    \begin{macrocode}
\LoadClass { ctexrep }
%    \end{macrocode}
%
% 载入各宏包。其中，\pkg{amsmath} 必须在 \pkg{unicode-math} 之前引入。
%    \begin{macrocode}
\RequirePackage
  {
    amsmath,
    unicode-math,
    geometry,
    fancyhdr,
    footmisc,
    graphicx,
    caption,
    bicaption,
    xcolor,
    xeCJKfntef
  }
%    \end{macrocode}
%
% \subsection{页面布局}
%
% 使用 \pkg{geometry} 宏包设置纸张大小、页面边距以及页眉高度。
%
% 根据《华东师范大学本科生毕业论文（设计）格式要求》的相关要求，
% 学士学位论文的页边距应设置为：上边距 \qty{2.5}{cm}，
% 下边距 \qty{2.0}{cm}，左边距 \qty{3.0}{cm}，右边距 
% \qty{2.5}{cm}。
%
% 根据《华东师范大学博士、硕士学位论文基本格式要求》的相关要求，
% 博士、硕士学位论文的页边距应按 Word 格式 A4 纸（“页面设置”按
% Word 默认值）设置。在 Word 提供的默认的页面设置中，左、右
% 边距均为 \qty{1.25}{in}；上、下边距均为 \qty{1}{in}。
%
% 在 Word 的默认设置中，页眉至顶端的距离设置为 \qty{15}{mm}，
% 页眉文字的字号为小五（即 \qty{9}{bp}），边框与正文的间距默认为 \qty{1}{bp}，。
% 在学士学位论文中，由于上边距为 \qty{25}{mm}，于是页眉与正文的间距应为
% \[ \qty{25}{mm}-\qty{15}{mm}-\qty{1}{bp}-1.296875\times\qty{9}{bp} = \qty{10}{mm}-\qty{12.6713575}{bp} \]
%
% 在博士、硕士学位论文中，由于上边距为 \qty{1}{in}，于是页眉与正文的间距应为
% \[ \qty{1}{in}-\qty{15}{mm}-\qty{1}{bp}-1.296875\times\qty{9}{bp} = \qty{1}{in}-\qty{15}{mm}-\qty{12.6713575}{bp} \]
%
% 在 Word 的默认设置中，页脚至底端的距离设置为 \qty{17.5}{mm}，
% 在学士学位论文中，由于下边距为 \qty{20}{mm}，若页脚至底端的距离仍为 \qty{17.5}{mm}，
% 则页脚高度仅为 \qty{2.5}{mm}，不足以容纳小五号字体。本模板将页脚与正文的间距设置
% 为 \qty{6}{mm}，这样页脚与正文之间将会留有一定的空白，也较为美观。
%
% 在博士、硕士学位论文中，由于下边距为 \qty{1}{in}，于是页脚与正文的间距应为 $\qty{1}{in}-\qty{17.5}{mm} $。
%
%    \begin{macrocode}
\bool_if:NTF \g_@@_is_bachelor_bool
  {
    \geometry
      {
        paper     = a4paper,
        top       = 2.5 cm,
        bottom    = 2.0 cm,
        left      = 3.0 cm,
        right     = 2.5 cm,
        headsep   = \dim_eval:n { 10 mm - 12.6713575 bp },
        footskip  = 6 mm
      }
  }
  {
    \geometry
      {
        paper     = a4paper,
        hmargin   = 1.25 in,
        vmargin   = 1 in,
        headsep   = \dim_eval:n { 1 in - 15 mm - 12.6713575 bp },
        footskip  = \dim_eval:n { 1 in - 17.5 mm }
      }
  }
%    \end{macrocode}
%
% \subsection{字体配置}
%
% \subsubsection{中西文字体配置}
%
% \begin{variable}{\g_@@_cjk_fontset_tl,\g_@@_fontset_tl}
% 存放字体选项值。
%    \begin{macrocode}
\tl_new:N \g_@@_cjk_fontset_tl
\tl_new:N \g_@@_fontset_tl
%    \end{macrocode}
% \end{variable}
%
% \begin{macro}{style/font}
% 定义西文字体选项。
%    \begin{macrocode}
\keys_define:nn { ecnu / style }
  {
    font  .choices:nn =
      { lm, times, xits }
      { \tl_set_eq:NN \g_@@_fontset_tl \l_keys_choice_tl },
    font  .value_required:n = true,
    font  .initial:n  = lm
  }
%    \end{macrocode}
% \end{macro}
%
% \begin{macro}{style/font-cjk}
% 定义中文字体选项。
%    \begin{macrocode}
\keys_define:nn { ecnu / style }
  {
    font-cjk  .choices:nn =
      { auto, fandol, mac, windows }
      {
        \tl_gset_eq:NN \g_@@_cjk_fontset_tl \l_keys_choice_tl
        \tl_if_eq:NnT \g_@@_cjk_fontset_tl { auto }
          {
            \sys_if_platform_windows:TF
              { \tl_gset:Nn \g_@@_cjk_fontset_tl { windows } }
              {
                \ctex_if_platform_macos:TF
                  { \tl_gset:Nn \g_@@_cjk_fontset_tl { mac    } }
                  { \tl_gset:Nn \g_@@_cjk_fontset_tl { fandol } }
              }
          }
      },
    font-cjk  .value_required:n = true,
    font-cjk  .initial:n  = auto
  }
%    \end{macrocode}
% \end{macro}
%
% \begin{macro}{
%   \@@_setmainfont:nn,
%   \@@_setsansfont:nn,
%   \@@_setmonofont:nn,
%   \@@_setmathfont:nn}
% 用于设置西文字体的辅助函数。代码来源于 \pkg{fontspec} 和 \pkg{unicode-math}。
% \begin{arguments}
%   \item 字体名
%   \item 选项
% \end{arguments}
%    \begin{macrocode}
\cs_new_protected:Npn \@@_setmainfont:nn #1#2
  { \__fontspec_main_setmainfont:nn {#2} {#1} }
\cs_new_protected:Npn \@@_setsansfont:nn #1#2
  { \__fontspec_main_setsansfont:nn {#2} {#1} }
\cs_new_protected:Npn \@@_setmonofont:nn #1#2
  { \__fontspec_main_setmonofont:nn {#2} {#1} }
\cs_new_protected:Npn \@@_setmathfont:nn #1#2
  { \__um_setmathfont:nn {#2} {#1} }
%    \end{macrocode}
% \end{macro}
%
% \begin{macro}{\@@_set_family:nnn,\@@_set_family:xnn,\@@_switch_family:n}
% 封装 CJK 字体族的设定和切换命令。
%    \begin{macrocode}
\cs_new_eq:NN \@@_set_family:nnn  \xeCJK_set_family:nnn
\cs_new_eq:NN \@@_switch_family:n \xeCJK_switch_family:n
\cs_generate_variant:Nn \@@_set_family:nnn { x }
%    \end{macrocode}
% \end{macro}
%
% \begin{macro}{
%   \@@_setCJKmainfont:nn,
%   \@@_setCJKsansfont:nn,
%   \@@_setCJKmonofont:nn}
% 用于设置中文字体的辅助函数。代码来源于 \pkg{xeCJK} 和 \pkg{ctex} 宏包。
%    \begin{macrocode}
\cs_new_protected:Npn \@@_setCJKmainfont:nn #1#2
  { \@@_set_family:xnn { \CJKrmdefault } {#2} {#1} }
\cs_new_protected:Npn \@@_setCJKsansfont:nn #1#2
  { \@@_set_family:xnn { \CJKsfdefault } {#2} {#1} }
\cs_new_protected:Npn \@@_setCJKmonofont:nn #1#2
  { \@@_set_family:xnn { \CJKttdefault } {#2} {#1} }
%    \end{macrocode}
% \end{macro}
%
% \begin{macro}{
%   \@@_set_cjk_font_kai:nn,
%   \ecnu@kai}
% 楷体需要单独设置。
%    \begin{macrocode}
\cs_new_protected:Npn \@@_set_cjk_font_kai:nn #1#2
  { \@@_set_family:nnn { ecnu@kai } {#2} {#1} }
\cs_new_protected:Npn \ecnu@kai
  { \@@_switch_family:n { ecnu@kai } }
%    \end{macrocode}
% \end{macro}
%
% \begin{macro}{
%   \@@_cjk_font_options:,
%   \@@_setCJKmainfont:n,
%   \@@_setCJKsansfont:n,
%   \@@_setCJKmonofont:n}
% 将 bold、italic 和 bold italic 统一按照 roman 设置。
%    \begin{macrocode}
\tl_const:Nn \@@_cjk_font_options:
 { UprightFont = *, ItalicFont = *, AutoFakeBold = true }
\cs_new_protected:Npx \@@_setCJKmainfont:n   #1
  { \@@_setCJKmainfont:nn   {#1} { \@@_cjk_font_options: } }
\cs_new_protected:Npx \@@_setCJKsansfont:n   #1
  { \@@_setCJKsansfont:nn   {#1} { \@@_cjk_font_options: } }
\cs_new_protected:Npx \@@_setCJKmonofont:n   #1
  { \@@_setCJKmonofont:nn   {#1} { \@@_cjk_font_options: } }
\cs_new_protected:Npx \@@_set_cjk_font_kai:n #1
  { \@@_set_cjk_font_kai:nn {#1} { \@@_cjk_font_options: } }
%    \end{macrocode}
% \end{macro}
%
% \begin{macro}{
%   \setmainfont,
%   \setsansfont,
%   \setmonofont,
%   \setmathfont,
%   \setCJKmainfont,
%   \setCJKsansfont,
%   \setCJKmonofont}
% 重新定义以上宏包提供的字体选择命令。
% 将其放在导言区末尾，使得用户配置不会被模板配置所覆盖。
%    \begin{macrocode}
\cs_new_protected:Npn \@@_set_font_helper:n #1
  {
    \exp_args:Nc \RenewDocumentCommand { set #1 font } { O { } m O { } }
      {
        \ctex_at_end_preamble:n
          { \use:c { @@_set #1 font:nn } {##2} { ##1, ##3 } }
      }
  }
\clist_map_inline:nn { main, sans, mono, math }
  { \@@_set_font_helper:n {#1} }
\clist_map_inline:nn { CJKmain, CJKsans, CJKmono }
  { \@@_set_font_helper:n {#1} }
%    \end{macrocode}
% \end{macro}
%
%
% \begin{macro}{\@@_load_font_times:}
% Times 系列西文字体。
%    \begin{macrocode}
\cs_new_protected:Npn \@@_load_font_times:
  {
    \@@_setmainfont:nn { Times~ New~ Roman    } { }
    \@@_setsansfont:nn { Arial                } { }
    \@@_setmonofont:nn { Courier~ New         } { }
    \@@_setmathfont:nn { XITSMath-Regular.otf }
      { BoldFont = XITSMath-Bold.otf }
  }
%    \end{macrocode}
% \end{macro}
%
% \begin{macro}{\@@_load_font_xits:}
% XITS 系列西文字体。
%    \begin{macrocode}
\cs_new_protected:Npn \@@_load_font_xits:
  {
    \@@_setmainfont:nn { XITS }
      {
        Extension      = .otf,
        UprightFont    = *-Regular,
        BoldFont       = *-Bold,
        ItalicFont     = *-Italic,
        BoldItalicFont = *-BoldItalic,
      }
    \@@_setsansfont:nn { texgyreheros }
      {
        Extension      = .otf,
        UprightFont    = *-regular,
        BoldFont       = *-bold,
        ItalicFont     = *-italic,
        BoldItalicFont = *-bolditalic
      }
    \@@_setmonofont:nn { texgyrecursor }
      {
        Extension      = .otf,
        UprightFont    = *-regular,
        BoldFont       = *-bold,
        ItalicFont     = *-italic,
        BoldItalicFont = *-bolditalic,
        Ligatures      = CommonOff
      }
    \@@_setmathfont:nn { XITSMath-Regular.otf }
      { BoldFont = XITSMath-Bold.otf }
  }
%    \end{macrocode}
% \end{macro}
%
% \begin{macro}{\@@_load_font_lm:}
% Latin Modern 系列西文字体。由于已作为默认字体，所以仅需额外处理数学部分。
%    \begin{macrocode}
\cs_new_protected:Npn \@@_load_font_lm:
  { \@@_setmathfont:nn { latinmodern-math.otf } { } }
%    \end{macrocode}
% \end{macro}
%
% \begin{macro}{\@@_load_cjk_font_fandol:}
% Fandol 字库中文字体。由于其安装在 TeX 发行版中，故使用文件名调用。
%    \begin{macrocode}
\cs_new_protected:Npn \@@_load_cjk_font_fandol:
  {
    \@@_setCJKmainfont:nn   { FandolSong }
      {
        Extension      = .otf,
        UprightFont    = *-Regular,
        BoldFont       = *-Bold,
        ItalicFont     = *-Regular,
        BoldItalicFont = *-Bold
      }
    \@@_setCJKsansfont:nn   { FandolHei  }
      {
        Extension      = .otf,
        UprightFont    = *-Regular,
        BoldFont       = *-Bold,
        ItalicFont     = *-Regular,
        BoldItalicFont = *-Bold
      }
    \@@_setCJKmonofont:nn   { FandolFang }
      {
        Extension      = .otf,
        UprightFont    = *-Regular,
        BoldFont       = *-Regular,
        ItalicFont     = *-Regular,
        BoldItalicFont = *-Regular
      }
    \@@_set_cjk_font_kai:nn { FandolKai  }
      {
        Extension      = .otf,
        UprightFont    = *-Regular,
        BoldFont       = *-Regular,
        ItalicFont     = *-Regular,
        BoldItalicFont = *-Regular
      }
  }
%    \end{macrocode}
% \end{macro}
%
% \begin{macro}{\@@_load_cjk_font_mac:}
% macOS 自带的中文字体。
%    \begin{macrocode}
\cs_new_protected:Npn \@@_load_cjk_font_mac:
  {
    \@@_setCJKmainfont:nn   { STSongti-SC }
      {
        UprightFont    = *-Light,
        BoldFont       = *-Bold,
        ItalicFont     = *-Light,
        BoldItalicFont = *-Bold
      }
    \@@_setCJKsansfont:nn   { STHeitiSC   }
      {
        UprightFont    = *-Medium,
        BoldFont       = *-Medium,
        ItalicFont     = *-Medium,
        BoldItalicFont = *-Medium
      }
    \@@_setCJKmonofont:n    { STFangsong  }
    \@@_set_cjk_font_kai:nn { STKaitiSC   }
      {
        UprightFont    = *-Regular,
        BoldFont       = *-Bold,
        ItalicFont     = *-Regular,
        BoldItalicFont = *-Bold
      }
  }
%    \end{macrocode}
% \end{macro}
%
% \begin{macro}{\@@_load_cjk_font_windows:}
% Windows 自带的中文字体。
%    \begin{macrocode}
\cs_new_protected:Npn \@@_load_cjk_font_windows:
  {
    \@@_setCJKmainfont:n   { SimSun   }
    \@@_setCJKsansfont:n   { SimHei   }
    \@@_setCJKmonofont:n   { FangSong }
    \@@_set_cjk_font_kai:n { KaiTi    }
  }
%    \end{macrocode}
% \end{macro}
%
% \begin{macro}{\@@_load_font:}
% 字体加载命令。
%    \begin{macrocode}
\cs_new_protected:Npn \@@_load_font:
  {
    \use:c { @@_load_font_ \g_@@_fontset_tl :         }
    \use:c { @@_load_cjk_font_ \g_@@_cjk_fontset_tl : }
  }
\ctex_at_end_preamble:n { \@@_load_font: }
%    \end{macrocode}
% \end{macro}
%
% \subsubsection{数学字体配置}
%
% 数学表达式中表示变量的拉丁字母和希腊字母均应当使用斜体。
%    \begin{macrocode}
\keys_set:nn { unicode-math }
  {
    math-style = ISO,
    bold-style = ISO,
  }
%    \end{macrocode}
%
% \subsubsection{其他配置}
%
% 声明 \cs{emph} 样式序列，使得中文下的强调命令用楷体显示。
%    \begin{macrocode}
\DeclareEmphSequence
  {
    \itshape \ecnu@kai,
    \upshape \CJKfamily { \CJKfamilydefault },
  }
%    \end{macrocode}
%
%
% \subsection{章节标题结构}
%
% 将标题层级设为五层。
%    \begin{macrocode}
\setcounter{secnumdepth}{5}
%    \end{macrocode}
%
% 设置本科生毕业论文章节的字体大小、缩进以及前后间距。本科生
% 毕业论文正文层次不超过 5 层次。由于文科模板与理科模板中
% 四级标题与五级标题的缩进不同，我们稍后再进行处理。
%    \begin{macrocode}
\bool_if:NTF \g_@@_is_bachelor_bool
  {
    \keys_set:nn { ctex }
      {
        chapter = {
          format        = \sffamily\normalsize,
          pagestyle     = fancy,
          beforeskip    = 1ex,
          afterskip     = 2.5ex plus .2ex,
          fixskip       = true,
          tocline       = {\CTEXifname{\protect\numberline{\CTEXthechapter}}{}#2}
        },
        section = {
          format        = \sffamily\normalsize,
          beforeskip    = 2.5ex plus 1ex minus .2ex,
          afterskip     = 2.5ex plus .2ex,
          fixskip       = true,
        },
        subsection = {
          format        = \sffamily\normalsize,
          beforeskip    = 2.5ex plus 1ex minus .2ex,
          afterskip     = 2.5ex plus .2ex,
          indent        = 1 em,
          fixskip       = true,
        },
        subsubsection   = {
          format        = \sffamily\normalsize,
          beforeskip    = 2.5ex plus 1ex minus .2ex,
          afterskip     = 2.5ex plus .2ex,
          indent        = 2 em,
          fixskip       = true,
        },
        paragraph = {
          format        = \sffamily\normalsize,
          beforeskip    = 2.5ex plus 1ex minus .2ex,
          afterskip     = 2.5ex plus .2ex,
          runin         = false,
          fixskip       = true,
        },
      }
  }
%    \end{macrocode}
%
% 设置博士、硕士学位论文章节的字体大小、缩进以及前后间距。
%    \begin{macrocode}
  {
    \keys_set:nn { ctex }
      {
        chapter = {
          format        = \zihao {3} \sffamily \bfseries \centering,
          pagestyle     = fancy,
          beforeskip    = 4ex,
          afterskip     = 6ex,
          fixskip       = true,
          tocline       = { \CTEXifname{\protect\numberline{\CTEXthechapter}}{}#2 }
        },
        section = {
          format        = \zihao {4} \bfseries \sffamily,
          beforeskip    = 3.5ex plus 1ex minus .2ex,
          afterskip     = 3.5ex plus .2ex,
          fixskip       = true,
          tocline       = { \CTEXnumberline {#1} #2 },
        },
        subsection = {
          format        = \bfseries \sffamily \normalsize,
          beforeskip    = 3.5ex plus 1ex minus .2ex,
          afterskip     = 3.5ex plus .2ex,
          fixskip       = true,
        },
        subsubsection   = {
          format        = \sffamily \normalsize,
          beforeskip    = 3.5ex plus 1ex minus .2ex,
          afterskip     = 3ex plus .2ex,
          indent        = 0em,
          fixskip       = true,
        },
        paragraph = {
          format        = \sffamily \normalsize,
          beforeskip    = 3.5ex plus 1ex minus .2ex,
          afterskip     = 3ex plus .2ex,
          indent        = 2em,
          runin         = false,
          fixskip       = true,
        },
      }
  }
%    \end{macrocode}
%
% \begin{variable}{\g_@@_numbering_tl}
% 论文章节编号样式类型。
%    \begin{macrocode}
\tl_new:N \g_@@_numbering_tl
%    \end{macrocode}
% \end{variable}
%
% \begin{macro}{style/numbering}
% 定义 \opt{style/numbering} 选项。
%    \begin{macrocode}
\keys_define:nn { ecnu / style }
  {
    numbering .choice:,
%    \end{macrocode}
% \end{macro}
%
% 设置理科模板的章节标题。理科模板的四级标题前的缩进为 2 个字符；
% 五级标题前的缩进为 3 个字符。
%    \begin{macrocode}
    numbering / arabic  .code:n = 
      {
        \tl_set:Nn \g_@@_numbering_tl { arabic }
        \bool_if:NTF \g_@@_is_bachelor_bool
          {
            \keys_set:nn { ctex }
              {
                chapter = {
                  name          = {,、},
                  number        = {\arabic{chapter}},
                  numberformat  = \rmfamily,
                  aftername     = ,
                },
                section = {
                  name          = {,},
                  number        = {\arabic{chapter}.\arabic{section}},
                  numberformat  = \rmfamily,
                },
                subsection = {
                  name          = {,},
                  number        = 
                    {\arabic{chapter}.\arabic{section}.\arabic{subsection}},
                  numberformat  = \rmfamily,
                }, 
                subsubsection = {
                  name          = {,},
                  number        = 
                    {
                      \arabic{chapter}.\arabic{section}.
                      \arabic{subsection}.\arabic{subsubsection}
                    },
                  numberformat  = \rmfamily,
                  indent        = 2 em,
                }, 
                paragraph = {
                  name          = {,）},
                  number        = {\arabic{paragraph}},
                  numberformat  = \rmfamily,
                  aftername     = ,
                  indent        = 3 em,
                }
              }
          }
          {
            \keys_set:nn { ctex }
              {
                chapter = {
                  name          = {第,章},
                  number        = {\chinese{chapter}},
                },
                section = {
                  name          = {,},
                  number        = {\arabic{chapter}.\arabic{section}},
                  numberformat  = \rmfamily,
                },
                subsection = {
                  name          = {,},
                  number        = 
                    {\arabic{chapter}.\arabic{section}.\arabic{subsection}},
                  numberformat  = \rmfamily,
                }, 
                subsubsection = {
                  name          = {,},
                  number        = 
                    {
                      \arabic{chapter}.\arabic{section}.
                      \arabic{subsection}.\arabic{subsubsection}
                    },
                  numberformat  = \rmfamily,
                }, 
                paragraph = {
                  name          = {,）},
                  number        = {\arabic{paragraph}},
                  numberformat  = \rmfamily,
                  aftername     = ,
                }
              }
          }
      },
%    \end{macrocode}
%
% 设置文科模板章节标题。文科模板的四级标题前的缩进为 2 个字符；
% 五级标题前的缩进为 4 个字符。
%    \begin{macrocode}
    numbering / chinese .code:n = 
      {
        \tl_set:Nn \g_@@_numbering_tl { chinese }
        \bool_if:NTF \g_@@_is_bachelor_bool
          {
            \keys_set:nn { ctex }
              {
                chapter = {
                  name          = {,、},
                  number        = {\chinese{chapter}},
                  numberformat  = \sffamily,
                },
                section = {
                  name          = {（,）},
                  number        = {\chinese{section}},
                  numberformat  = \sffamily,
                  aftername     = ,
                }, 
                subsection = {
                  name          = {,.},
                  number        = {\arabic{subsection}},
                  indent        = 2 em,
                }, 
                subsubsection = {
                  name          = {（,）},
                  number        = {\arabic{subsubsection}},
                  aftername     = {\ },
                  indent        = 2 em,
                }, 
                paragraph = {
                  name          = {,）},
                  number        = {\arabic{paragraph}},
                  aftername     = ,
                  indent        = 4 em,
                }
              }
          }
          {
            \keys_set:nn { ctex }
              {
                chapter = {
                  name          = {第,章},
                  number        = {\chinese{chapter}},
                  numberformat  = \sffamily,
                },
                section = {
                  name          = {第,节},
                  number        = {\chinese{section}},
                  numberformat  = \sffamily,
                }, 
                subsection = {
                  name          = {,、},
                  number        = {\chinese{subsection}},
                  numberformat  = \sffamily,
                  aftername     = ,
                }, 
                subsubsection = {
                  name          = {（,）},
                  number        = {\chinese{subsubsection}},
                  numberformat  = \sffamily,
                  aftername     = ,
                }, 
                paragraph = {
                  name          = {,）},
                  number        = {\arabic{paragraph}},
                  numberformat  = \rmfamily,
                  aftername     = ,
                }
              }
          }
        },
%    \end{macrocode}
%
% 处理未知选项。
%    \begin{macrocode}
    numbering / unknown   .code:n = { \@@_error:n { unknown-value } },
    numbering .value_required:n = true,
    numbering .initial:n = arabic
  }
\@@_msg_new:nn { unknown-value }
  { The~ value~ "\l_keys_value_tl"~ is~ unknown. }
%    \end{macrocode}
%
% \begin{macro}{\@@_title:n,\@@_title:V}
% 手动生成章的标题，用于摘要、参考文献等。
% 在本科生毕业论文中，这些标题的字体为黑体加粗、字号为小三、对齐方式为居中；
% 在硕士、博士学位论文中，这些标题的字体为黑体加粗、字号为三号、对齐方式为居中。
%    \begin{macrocode}
\cs_new_protected:Npn \@@_title:n #1
  {
    \group_begin:
      \ctexset 
        { 
          chapter / numbering = false,
          chapter / format    = 
            {
              \sffamily\bfseries\centering
              \bool_if:NTF \g_@@_is_bachelor_bool
                { \zihao {-3} } { \zihao {3} }
            },
        }
      \chapter {#1}
    \group_end:
  }
\cs_generate_variant:Nn \@@_title:n { V }
%    \end{macrocode}
% \end{macro}
%
% \begin{macro}{\@@_abs_title:n,\@@_abs_title:V}
% 手动生成摘要页的中文标题。这些标题的字体大小为小三。
%    \begin{macrocode}
\cs_new_protected:Npn \@@_abs_title:n #1
  {
    \group_begin:
      \ctexset 
        { 
          chapter / numbering   = false,
          chapter / format      = \sffamily\bfseries\centering\zihao{-3},
        }
      \chapter* {#1}
    \group_end:
  }
\cs_generate_variant:Nn \@@_abs_title:n { V }
%    \end{macrocode}
% \end{macro}
%
% \begin{macro}{\@@_absEN_title:n,\@@_absEN_title:V}
% 手动生成摘要页的英文标题。这些标题的字体大小为小三。
%    \begin{macrocode}
\cs_new_protected:Npn \@@_absEN_title:n #1
  {
    \group_begin:
      \ctexset
        {
          chapter / numbering   = false,
          chapter / format      = \rmfamily\bfseries\centering\zihao{-3},
        }
      \chapter* {#1}
    \group_end:
  }
\cs_generate_variant:Nn \@@_absEN_title:n { V }
%    \end{macrocode}
% \end{macro}
%
% \subsection{页眉页脚设置}
%
% 清除默认页眉页脚格式。
%    \begin{macrocode}
\fancyhf { }
%    \end{macrocode}
%
% \begin{macro}{\@@_pg_header_content:}
% 博士、硕士学位论文的页眉内容。
% 
% 根据《华东师范大学博士、硕士学位论文基本格式要求》的相关要求，
% 博士、硕士学位论文的页眉应写明“华东师范大学博（硕）士学位论文”字样；
% 专业学位论文应写明“华东师范大学博（硕）士专业学位论文”字样。
%    \begin{macrocode}
\cs_new_protected:Npn \@@_pg_header_content:
  {
    \prop_get:NVN \c_@@_thesis_type_prop
      \g_@@_thesis_type_tl \l_@@_tmpa_tl
    \c_@@_university_name_tl \l_@@_tmpa_tl
    \tl_if_eq:NnT \l_@@_info_degree_type_tl
      { professional } { \c_@@_professional_tl }
    \c_@@_degree_tl \c_@@_thesis_tl
  }
%    \end{macrocode}
% \end{macro}
%
% 本科生毕业论文的页眉样式和页脚样式。
% 
% 根据《本科生毕业论文（设计）模板》的相关要求，
% 页眉的内容为毕业论文（设计）题目，格式为黑体、小五号；
% 页脚的内容为阿拉伯数字页码，格式为小五、居中。
%    \begin{macrocode}
\bool_if:NTF \g_@@_is_bachelor_bool
  {
    \fancyhead[C]
      {
        \footnotesize \sffamily
        \l_@@_info_title_trim_tl
      }
    \fancyfoot[C] { \footnotesize \thepage }
  }
%    \end{macrocode}
%
% 博士、硕士学位论文的页眉样式和页脚样式。
% 
% 页眉按照 Word 默认值排版。
% 双面复印时页脚的页码分别按左右侧排列。
%    \begin{macrocode}
  {
    \fancyhead[C] { \footnotesize \@@_pg_header_content: }
    \bool_if:NTF \g_@@_twoside_bool
      { \fancyfoot[LE,RO] { \footnotesize \thepage } }
      { \fancyfoot[C] { \footnotesize \thepage } }
  }
%    \end{macrocode}
%
%
% 将页眉分隔线的宽度设置为 \qty{0.5}{bp} 。
%    \begin{macrocode}
\RenewDocumentCommand \headrulewidth { } { 0.5 bp }
%    \end{macrocode}
%
% \begin{macro}{\cleardoublepage}
% 重定义 \tn{cleardoublepage}，使得偶数页面在没有内容时也不显示页眉页脚，见
%    \begin{macrocode}
\RenewDocumentCommand \cleardoublepage { }
  {
    \clearpage
    \bool_if:NT \g_@@_twoside_bool
      {
        \int_if_odd:nF \c@page
          { \hbox:n { } \thispagestyle { empty } \newpage }
      }
  }
%    \end{macrocode}
% \end{macro}
%
% 由于 \pkg{ctex} 宏包使用 \opt{heading} 选项后，会将页面格式设置为 |headings|。
% 因此需要在调用 \pkg{ctex} 后重新将 \opt{pagestyle} 设置为 |fancy|。
%    \begin{macrocode}
\pagestyle { fancy }
%    \end{macrocode}
%
% \subsection{脚注}
%
% 各种脚注编号样式的名称。
%    \begin{macrocode}
\cs_new_protected:Npn \@@_define_fn_style:nn #1#2
  { \tl_const:cn { c_@@_fn_style_ #1 _tl } {#2} }
\clist_map_inline:nn
  {
    { plain       } { plain       },
    { xits        } { xits        },
    { libertinus  } { libertinus  },
  }
  { \@@_define_fn_style:nn #1 }
%    \end{macrocode}
%
% \begin{variable}{\l_@@_fn_style_tl}
% 保存当前使用的脚注编号样式。
%    \begin{macrocode}
\tl_new:N \l_@@_fn_style_tl
%    \end{macrocode}
% \end{variable}
%
% 提供设置脚注选项的接口。
%    \begin{macrocode}
\keys_define:nn { ecnu / style }
  {
%    \end{macrocode}
%
%    \begin{macrocode}
    footnote  .choices:nn = { plain, xits, libertinus }
%    \end{macrocode}
%
%    \begin{macrocode}
      {
        \tl_set_eq:NN \l_@@_fn_style_tl \l_keys_choice_tl
      },
    footnote  .value_required:n = true,
    footnote  .initial:n = plain
  }
%    \end{macrocode}
%
% \begin{macro}{\@@_fn_symbol_libertinus:n}
% 使用 libertinus 字体作为编号字体。在该字体下，1至 20 为数字，
% 21 至 46 为小写英文字母，47 至 72 为大写英文字母。
%    \begin{macrocode}
\cs_new:Npn \@@_fn_symbol_libertinus:n #1
  {
    \int_compare:nTF { #1 >= 21 }
      {
        \int_compare:nTF { #1 >= 47 }
          { \@@_symbol:n { \int_eval:n { "24B6 - 47 + #1 } } }
          { \@@_symbol:n { \int_eval:n { "24D0 - 21 + #1 } } }
      }
      { \@@_symbol:n { \int_eval:n { "2460 - 1 + #1 } } }
  }
%    \end{macrocode}
% \end{macro}
%
% \begin{macro}{\@@_fn_symbol_xits:n}
% 使用 xits 字体作为编号字体。在该字体下，1 至 9 为数字，10 至 35
% 为小写英文字母，36 至 61 为大写英文字母。
%    \begin{macrocode}
\cs_new:Npn \@@_fn_symbol_xits:n #1
  {
    \int_compare:nTF { #1 >= 10 }
      {
        \int_compare:nTF { #1 >= 36 }
          { \@@_symbol:n { \int_eval:n { "24B6 - 36 + #1 } } }
          { \@@_symbol:n { \int_eval:n { "24D0 - 10 + #1 } } }
      }
      { \@@_symbol:n { \int_eval:n { "2460 - 1 + #1 } } }
  }
%    \end{macrocode}
% \end{macro}
%
% \begin{macro}{\thefootnote}
% 重定义脚注编号。
%    \begin{macrocode}
\cs_set:Npn \thefootnote { \@@_footnote_number:N \c@footnote }
%    \end{macrocode}
% \end{macro}
%
% \begin{macro}{\@@_footnote_number:N}
% 脚注编号样式。
%    \begin{macrocode}
\cs_new:Npn \@@_footnote_number:N #1
  {
    \tl_case:NnF \l_@@_fn_style_tl
      {
        \c_@@_fn_style_plain_tl
          { \int_use:N #1 }
        \c_@@_fn_style_xits_tl
          {
            \fontspec { XITS-Regular.otf }
            \@@_fn_symbol_xits:n { #1 }
          }
        \c_@@_fn_style_libertinus_tl
          {
            \fontspec { LibertinusSerif-Regular.otf }
            \@@_fn_symbol_libertinus:n { #1 }
          }
      }
      { \int_use:N #1 }
  }
%    \end{macrocode}
% \end{macro}
%
% \begin{macro}{\@makefntext}
% 在本科生毕业论文中，脚注没有缩进。故重定义内部脚注文字命令，
% 使生成的脚注无缩进。
%    \begin{macrocode}
\cs_set:Npn \@makefntext #1
  {
    \dim_set:Nn \parindent
      { \bool_if:NTF \g_@@_is_bachelor_bool { 0 em } { 2 em } }
    \mode_leave_vertical:
    \textsuperscript{\@thefnmark} #1
  }
%    \end{macrocode}
% \end{macro}
%
% \subsection{图表绘制；浮动体}
%
% 设置浮动体 figure 和 table 的标题样式。
%    \begin{macrocode}
\captionsetup
  {
    font          = small,
    labelsep      = quad,
    justification = centering
  }
%    \end{macrocode}
%
% 设置外文标题的名称。
%    \begin{macrocode}
\captionsetup [ figure ] [ bi-second ] { name=Figure }
\captionsetup [ table  ] [ bi-second ] { name=Table  }
%    \end{macrocode}
% 
% \begin{macro}{\@makefntext}
% 在本科生毕业论文中，图表的编号格式需要重新定义。
%    \begin{macrocode}
\bool_if:NT \g_@@_is_bachelor_bool
  {
    \cs_set:Npn \thefigure
      {  \thechapter - \int_to_arabic:n { \c@figure } }
    \cs_set:Npn \thetable
      {  \thechapter - \int_to_arabic:n { \c@table  } }
  }
%    \end{macrocode}
% \end{macro}
%
% \subsection{封面}
%
% \subsubsection{信息录入}
%
% \begin{variable}{\l_@@_info_title_tl,
%   \l_@@_info_title_en_tl,
%   \l_@@_info_author_tl,
%   \l_@@_info_author_en_tl,
%   \l_@@_info_student_id_tl,
%   \l_@@_info_department_tl,
%   \l_@@_info_department_en_tl,
%   \l_@@_info_major_tl,
%   \l_@@_info_major_en_tl,
%   \l_@@_info_supervisor_tl,
%   \l_@@_info_supervisor_en_tl,
%   \l_@@_info_academic_title_tl,
%   \l_@@_info_keywords_clist,
%   \l_@@_info_keywords_en_clist}
% 封面所需的一些字段（中外文标题、作者姓名、学号、院系、专业、
% 导师姓名、导师职称、中外文关键词）。
%    \begin{macrocode}
\clist_map_inline:nn
  {
    title, title_en,
    author, author_en,
    student_id,
    department, department_en,
    major, major_en,
    supervisor, supervisor_en,
    field, field_en,
    academic_title
  }
  { \tl_new:c { l_@@_info_ #1 _tl } }
\clist_new:N \l_@@_info_keywords_clist
\clist_new:N \l_@@_info_keywords_en_clist
%    \end{macrocode}
% \end{variable}
%
% \begin{variable}{\l_@@_info_degree_type_tl}
% 博士、硕士学位论文中的学位类型。
%    \begin{macrocode}
\tl_new:N \l_@@_info_degree_type_tl
%    \end{macrocode}
% \end{variable}
%
% \begin{variable}{\l_@@_info_cic_index_tl}
% 博士、硕士学位论文中的中图分类号。
%    \begin{macrocode}
\tl_new:N \l_@@_info_cic_index_tl
%    \end{macrocode}
% \end{variable}
%
% \begin{variable}{\l_@@_info_secret_level_tl}
% 博士、硕士学位论文中的密级。
%    \begin{macrocode}
\tl_new:N \l_@@_info_secret_level_tl
%    \end{macrocode}
% \end{variable}
%
% \begin{variable}{\l_@@_info_year_int,
%   \l_@@_info_month_int,
%   \l_@@_info_graduation_year_int}
% 论文完成年份、月份与毕业年份，以及博士、硕士学位论文中的论文完成日期。
%    \begin{macrocode}
\int_new:N   \l_@@_info_year_int
\int_new:N   \l_@@_info_month_int
\int_new:N   \l_@@_info_graduation_year_int
%    \end{macrocode}
% \end{variable}
%
% \begin{variable}{\l_@@_info_title_trim_tl,\l_@@_info_title_en_trim_tl}
% 去除空格和换行后的论文标题。
%    \begin{macrocode}
\tl_new:N \l_@@_info_title_trim_tl
\tl_new:N \l_@@_info_title_en_trim_tl
%    \end{macrocode}
% \end{variable}
%
% \begin{macro}{\@@_set_title:nn}
% 设置论文标题，并删除其中的换行和汉字间的空格。
%    \begin{macrocode}
\cs_generate_variant:Nn \regex_replace_all:nnN { nnc }
\cs_new_protected:Npn \@@_set_title:nn #1#2
  {
    \tl_set:cn { l_@@_info_ #1 _tl      } {#2}
    \tl_set:cn { l_@@_info_ #1 _trim_tl } {#2}
    \tl_replace_all:cnn { l_@@_info_ #1 _trim_tl } { \\ } { ~ }
    \regex_replace_all:nnc
      { \s+ } { \ }
      { l_@@_info_ #1 _trim_tl }
    \regex_replace_all:nnc
      { ([\x{4e00}-\x{9fff}]) \s+ ([\x{4e00}-\x{9fff}]) } { \1 \2 }
      { l_@@_info_ #1 _trim_tl }
  }
%    \end{macrocode}
% \end{macro}
%
% 定义 \opt{ecnu/info} 键值类。
%    \begin{macrocode}
\keys_define:nn { ecnu / info }
  {
%    \end{macrocode}
%
% \begin{macro}{info/degree}
% 博士、硕士学位论文的学位类型。
%    \begin{macrocode}
    degree          .choices:nn   =
      { academic, professional }
      { \tl_set_eq:NN \l_@@_info_degree_type_tl \l_keys_choice_tl },
    degree          .value_required:n = true,
    degree          .initial:n = academic,
%    \end{macrocode}
% \end{macro}
%
% \begin{macro}{info/clc}
% 博士、硕士学位论文的中图分类号。
%    \begin{macrocode}
    clc             .tl_set:N     = \l_@@_info_cic_index_tl,
%    \end{macrocode}
% \end{macro}
%
% \begin{macro}{info/secret-level}
% 博士、硕士学位论文的密级。
%    \begin{macrocode}
    secret-level    .tl_set:N     = \l_@@_info_secret_level_tl,
%    \end{macrocode}
% \end{macro}
%
% \begin{macro}{info/title,info/title-en}
% 论文的中外文标题。
%    \begin{macrocode}
    title           .code:n       = { \@@_set_title:nn { title    } {#1} },
    title-en        .code:n       = { \@@_set_title:nn { title_en } {#1} },
%    \end{macrocode}
% \end{macro}
%
% \begin{macro}{info/author,info/author-en}
% 论文作者姓名。
%    \begin{macrocode}
    author          .tl_set:N     = \l_@@_info_author_tl,
    author-en       .tl_set:N     = \l_@@_info_author_en_tl,
%    \end{macrocode}
% \end{macro}
%
% \begin{macro}{info/student-id}
% 论文作者学号。
%    \begin{macrocode}
    student-id      .tl_set:N     = \l_@@_info_student_id_tl,
%    \end{macrocode}
% \end{macro}
%
% \begin{macro}{info/department,info/department-en}
% 论文作者所在院系。
%    \begin{macrocode}
    department      .tl_set:N     = \l_@@_info_department_tl,
    department-en   .tl_set:N     = \l_@@_info_department_en_tl,
%    \end{macrocode}
% \end{macro}
%
% \begin{macro}{info/major,info/major-en}
% 论文作者就读专业。
%    \begin{macrocode}
    major           .tl_set:N     = \l_@@_info_major_tl,
    major-en        .tl_set:N     = \l_@@_info_major_en_tl,
%    \end{macrocode}
% \end{macro}
%
% \begin{macro}{info/field,info/field-en}
% 博士、硕士学位论文的研究方向或专业学位领域。
%    \begin{macrocode}
    field           .tl_set:N     = \l_@@_info_field_tl,
    field-en        .tl_set:N     = \l_@@_info_field_en_tl,
%    \end{macrocode}
% \end{macro}
%
% \begin{macro}{info/supervisor,info/supervisor-en}
% 论文指导教师姓名。
%    \begin{macrocode}
    supervisor      .tl_set:N     = \l_@@_info_supervisor_tl,
    supervisor-en   .tl_set:N     = \l_@@_info_supervisor_en_tl,
%    \end{macrocode}
% \end{macro}
%
% \begin{macro}{info/academic-title}
% 本科生毕业论文的指导教师职称。
%    \begin{macrocode}
    academic-title   .tl_set:N     = \l_@@_info_academic_title_tl,
%    \end{macrocode}
% \end{macro}
%
% \begin{macro}{info/keywords,info/keywords-en}
% 论文中外文关键词。
%    \begin{macrocode}
    keywords        .clist_set:N  = \l_@@_info_keywords_clist,
    keywords-en     .clist_set:N  = \l_@@_info_keywords_en_clist,
%    \end{macrocode}
% \end{macro}
%
% \begin{macro}{info/year, info/month}
% 论文完成年份与月份。
%    \begin{macrocode}
    year            .int_set:N    = \l_@@_info_year_int,
    month           .int_set:N    = \l_@@_info_month_int,
%    \end{macrocode}
% \end{macro}
%
% \begin{macro}{info/graduation-year}
% 本科生毕业论文的作者毕业年份。
%    \begin{macrocode}
    graduation-year .int_set:N    = \l_@@_info_graduation_year_int,
  }
%    \end{macrocode}
% \end{macro}
%
% \begin{macro}{\theyear}
% 记录论文完成年份。默认为当前年份。
%    \begin{macrocode}
\NewDocumentCommand \theyear { } 
  {
    \int_compare:nNnTF { \l_@@_info_year_int } = { \c_zero_int } 
      { \int_set_eq:NN \l_@@_info_year_int \c_sys_year_int } 
      { }
    \int_use:N \l_@@_info_year_int
  }
%    \end{macrocode}
% \end{macro}
%
% \begin{macro}{\themonth}
% 记录论文完成月份。默认为当前月份。
%    \begin{macrocode}
\NewDocumentCommand \themonth { } 
  {
    \int_compare:nNnTF { \l_@@_info_month_int } = { \c_zero_int } 
      { \int_set_eq:NN \l_@@_info_month_int \c_sys_month_int }
      { }
    \int_use:N \l_@@_info_month_int
  }
%    \end{macrocode}
% \end{macro}
%
% \begin{macro}{\graduationyear}
% 记录论文作者毕业届数。若 \opt{info/graduation-year} 未设置，则使用
% \opt{info/year} 的值；若\opt{info/year} 的值亦未设置，则使用
% 当前年份。
%    \begin{macrocode}
\NewDocumentCommand \graduationyear { } 
  {
    \int_compare:nNnTF { \l_@@_info_graduation_year_int } = { \c_zero_int } 
      {
        \int_compare:nNnTF { \l_@@_info_year_int } = { \c_zero_int } 
          {
            \int_set_eq:NN  \l_@@_info_graduation_year_int 
                            \c_sys_year_int
          } 
          {
            \int_set_eq:NN  \l_@@_info_graduation_year_int 
                            \l_@@_info_year_int
          }
      }
      { }
    \int_use:N  \l_@@_info_graduation_year_int
  }
%    \end{macrocode}
% \end{macro}
%
% \subsubsection{定义内部函数}
%
% \begin{macro}{\@@_spread_box:nn,\@@_spread_box:no}
% 分散对齐的水平盒子。
% \begin{arguments}
%   \item 宽度
%   \item 内容
% \end{arguments}
% 利用 \cs{tl_map_inline:nn} 在字符间插入 |hfil|；紧随其后的 \tn{unskip}
% 将会去掉最后一个 |hfil|。
%    \begin{macrocode}
\cs_new_protected:Npn \@@_spread_box:nn #1#2
  {
    \mode_leave_vertical:
    \hbox_to_wd:nn {#1} { \tl_map_inline:nn {#2} { ##1 \hfil } \unskip }
  }
\cs_generate_variant:Nn \@@_spread_box:nn { no }
%    \end{macrocode}
% \end{macro}
%
% \begin{macro}{\@@_center_box:nn,\@@_center_box:Vn}
% 居中对齐的水平盒子。
%    \begin{macrocode}
\cs_new_protected:Npn \@@_center_box:nn #1#2
  {
    \mode_leave_vertical:
    \hbox_to_wd:nn {#1} { \hfil #2 \hfil }
  }
\cs_generate_variant:Nn \@@_center_box:nn { Vn }
%    \end{macrocode}
% \end{macro}
%
% \begin{macro}{\@@_right_box:nn,\@@_right_box:no}
% 右对齐的水平盒子。
%    \begin{macrocode}
\cs_new_protected:Npn \@@_right_box:nn #1#2
  {
    \mode_leave_vertical:
    \hbox_to_wd:nn {#1} { \hfil #2 }
  }
\cs_generate_variant:Nn \@@_right_box:nn { no }
%    \end{macrocode}
% \end{macro}
%
% \begin{macro}{\@@_fixed_width_box:nn}
% 限宽盒子（允许换行）。
%    \begin{macrocode}
\cs_new:Npn \@@_fixed_width_box:nn #1#2
  { \parbox {#1} {#2} }
%    \end{macrocode}
% \end{macro}
%
% \begin{macro}{\@@_fixed_width_top_align_box:nn}
% 顶端对齐的限宽盒子（允许换行）。
%    \begin{macrocode}
\cs_new:Npn \@@_fixed_width_align_top_box:nn #1#2
  { \parbox [t] {#1} {#2} }
%    \end{macrocode}
% \end{macro}
%
% \begin{macro}{\@@_fixed_width_center_box:nn}
% 居中对齐的限宽盒子（允许换行）。
%    \begin{macrocode}
\cs_new:Npn \@@_fixed_width_center_box:nn #1#2
  { \parbox {#1} { \centering #2 } }
%    \end{macrocode}
% \end{macro}
%
% \begin{macro}{\@@_get_text_width:Nn,\@@_get_text_width:NV}
% 获取文本宽度，并存入 |dim| 型变量。
% \begin{arguments}
%   \item |dim| 型变量
%   \item 内容
% \end{arguments}
%    \begin{macrocode}
\cs_new:Npn \@@_get_text_width:Nn #1#2
  {
    \hbox_set:Nn \l_@@_tmpa_box {#2}
    \dim_set:Nn #1 { \box_wd:N \l_@@_tmpa_box }
  }
\cs_generate_variant:Nn \@@_get_text_width:Nn { NV }
%    \end{macrocode}
% \end{macro}
%
% \begin{macro}{\@@_line_spread:N,\@@_line_spread:n}
% 设置行距。
%    \begin{macrocode}
\cs_new:Npn \@@_line_spread:N #1
  { \linespread { \fp_use:N #1 } \selectfont }
\cs_new:Npn \@@_line_spread:n #1
  { \linespread {#1} \selectfont }
%    \end{macrocode}
% \end{macro}
%
% \subsubsection{封面各部件}
%
% \begin{macro}{\ULthickness}
% 设置 \pkg{xeCJKfntef} 宏包所依赖的 \pkg{ulem} 宏包的下划线宽度。
% 在本科生毕业论文中，下划线宽度一般为 \qty{1}{pt}；在博士、硕士学位
% 论文中，下划线宽度一般为 \qty{0.5}{pt} 。
%    \begin{macrocode}
\bool_if:NF \g_@@_is_bachelor_bool
  { \RenewDocumentCommand \ULthickness { } { 0.5 pt } }
%    \end{macrocode}
% \end{macro}
%
% \begin{macro}{\@@_cover_ug_thesis_info:,\@@_cover_pg_thesis_info:,\@@_cover_pg_thesis_info_en:}
% 封面左上角首行的论文类型信息。
%    \begin{macrocode}
\cs_new_protected:Npn \@@_cover_ug_thesis_info:
  {
    \prop_get:NVN \c_@@_thesis_type_prop
      \g_@@_thesis_type_tl \l_@@_tmpa_tl
    \graduationyear \c_@@_graduation_year_tl
    \c_@@_undergraduate_tl \l_@@_tmpa_tl
    \c_@@_degree_tl \c_@@_thesis_tl
  }
\cs_new_protected:Npn \@@_cover_pg_thesis_info:
  {
    \prop_get:NVN \c_@@_thesis_type_prop
      \g_@@_thesis_type_tl \l_@@_tmpa_tl
    \graduationyear \c_@@_graduation_year_tl
    \tl_if_eq:NnTF \l_@@_info_degree_type_tl { professional }
      {
        \l_@@_tmpa_tl \c_@@_professional_tl \c_@@_degree_tl
        \c_@@_postgraduate_tl 
      }
      { \c_@@_postgraduate_tl \l_@@_tmpa_tl }
    \c_@@_degree_tl \c_@@_thesis_tl
  }
\cs_new_protected:Npn \@@_cover_pg_thesis_info_en:
  {
    \tl_if_eq:NnT \g_@@_thesis_type_tl { master }
      {
        \tl_if_eq:NnTF \l_@@_info_degree_type_tl { professional }
          { \c_@@_thesis_en_tl } { \c_@@_dissertation_en_tl }
        \c_space_token \c_@@_for_en_tl \c_space_token
        \c_@@_masters_en_tl \c_space_token \c_@@_degree_en_tl
        \tl_if_eq:NnT \l_@@_info_degree_type_tl { professional }
          { \c_space_token ( \c_@@_professional_en_tl ) }
      }
    \tl_if_eq:NnT \g_@@_thesis_type_tl { doctor }
      {
        \c_@@_dissertation_en_tl \c_space_token
        \c_@@_for_en_tl \c_space_token
        \c_@@_doctoral_en_tl \c_space_token \c_@@_degree_en_tl
        \tl_if_eq:NnT \l_@@_info_degree_type_tl { professional }
          { \c_space_token ( \c_@@_professional_en_tl ) }
      }
    \c_space_token in \c_space_token \graduationyear
  }
%    \end{macrocode}
% \end{macro}
%
% \begin{macro}{\@@_cover_ug_school_num:,\@@_cover_pg_school_num:,\@@_cover_pg_school_num_en:}
% 封面顶部的学校代码信息。
%    \begin{macrocode}
\cs_new_protected:Npn \@@_cover_ug_school_num:
  {
    \c_@@_university_code_tl \c_@@_colon_tl
    \underline{\c_@@_university_code_val_tl}
  }
\cs_new_protected:Npn \@@_cover_pg_school_num:
  {
    \@@_spread_box:no { 4 em } { \c_@@_university_code_tl }
    \c_@@_colon_tl
    \uline{\@@_center_box:nn { 3 cm } { \c_@@_university_code_val_tl }}
  }
\cs_new_protected:Npn \@@_cover_pg_school_num_en:
  {
    \@@_get_text_width:Nn \l_@@_tmpa_dim
      {
        \zihao {5}
        \c_@@_university_code_en_tl \c_@@_colon_en_tl \c_space_token
      }
    \@@_fixed_width_box:nn { \l_@@_tmpa_dim }
      { \c_@@_university_code_en_tl \c_@@_colon_en_tl \c_space_token }
    \uline{\@@_center_box:nn { 3 cm } { \c_@@_university_code_val_tl }}
  }
%    \end{macrocode}
% \end{macro}
%
% \begin{macro}{\@@_cover_pg_student_id:,\@@_cover_pg_student_id_en:}
% 博士、硕士学位论文中封面顶部的学号信息。
%    \begin{macrocode}
\cs_new_protected:Npn \@@_cover_pg_student_id:
  {
    \@@_spread_box:no { 4 em } { \c_@@_info_student_id_tl }
    \c_@@_colon_tl
    \uline{ \@@_center_box:nn { 3 cm } { \l_@@_info_student_id_tl } }
  }
\cs_new_protected:Npn \@@_cover_pg_student_id_en:
  {
    \@@_get_text_width:Nn \l_@@_tmpa_dim
      {
        \zihao {5} \c_@@_university_code_en_tl
        \c_@@_colon_en_tl \c_space_token
      }
    \@@_fixed_width_box:nn { \l_@@_tmpa_dim }
      { \c_@@_info_student_id_en_tl \c_@@_colon_en_tl \c_space_token }
    \uline{ \@@_center_box:nn { 3 cm } { \l_@@_info_student_id_tl } }
  }
%    \end{macrocode}
% \end{macro}
%
% \begin{macro}{\@@_cover_pg_cic_index:}
% 博士、硕士学位论文中封面顶部的中图分类号信息。
%    \begin{macrocode}
\cs_new_protected:Npn \@@_cover_pg_cic_index:
  {
    \@@_spread_box:no { 4 em } { \c_@@_cic_index_tl } \c_@@_colon_tl
    \uline{ \@@_center_box:nn { 3 cm } { \l_@@_info_cic_index_tl } }
  }
%    \end{macrocode}
% \end{macro}
%
% \begin{macro}{\@@_cover_pg_secret_level:}
% 博士、硕士学位论文中封面顶部的密级信息。
%    \begin{macrocode}
\cs_new_protected:Npn \@@_cover_pg_secret_level:
  {
    \@@_spread_box:no { 4 em } { \c_@@_secret_level_tl  } \c_@@_colon_tl
    \uline{ \@@_center_box:nn { 3 cm } { \l_@@_info_secret_level_tl } }
  }
%    \end{macrocode}
% \end{macro}
%
% \begin{macro}{\@@_cover_ug_top:,\@@_cover_pg_top:,\@@_cover_pg_top_en:}
% 封面顶部信息部件。
%    \begin{macrocode}
\cs_new_protected:Npn \@@_cover_ug_top:
  {
    \@@_fixed_width_box:nn { \textwidth }
      {
        \bfseries \zihao { 4 }
        \@@_cover_ug_thesis_info: \hfill \@@_cover_ug_school_num:
      }
  }
\cs_new_protected:Npn \@@_cover_pg_top:
  {
    \@@_fixed_width_box:nn { \textwidth }
      {
        \@@_cover_pg_thesis_info:
        \skip_vertical:n { 0 em }
        \@@_cover_pg_cic_index: \hfill \@@_cover_pg_school_num:
        \skip_vertical:n { 0 em }
        \@@_cover_pg_secret_level: \hfill \@@_cover_pg_student_id:
      }
  }
\cs_new_protected:Npn \@@_cover_pg_top_en:
  {
    \@@_fixed_width_box:nn { \textwidth }
      {
        \@@_cover_pg_thesis_info_en:
        \skip_vertical:n { 0 em }
        \hbox:n {} \mode_leave_vertical:
        \hfill \@@_cover_pg_school_num_en:
        \skip_vertical:n { 0 em }
        \hfill \@@_cover_pg_student_id_en:
      }
  }
%    \end{macrocode}
% \end{macro}
%
% \begin{macro}{\@@_pg_title_fixed_ht_box:nnn}
% 博士、硕士学位论文封面中外文标题盒子（带下划线）。
%
% 由于 \LaTeXe{} 没有提供一个对长文本进行自动换行后为每一行加上下划线且居中的实现，故
% 此处尝试使用 \TeX{} 提供的 \cs{vsplit} 原语，通过手动换行来实现这一效果。
%    \begin{macrocode}
\cs_new_protected:Npn \@@_pg_title_fixed_ht_box:nnn #1#2#3
  {
    \vbox_top:n
      {
%    \end{macrocode}
%
% 先将标题内容放入一个竖直盒子中进行预排版。值得注意的是，此处需设置
% \cs{splittopskip} 使 \TeX{} 在 \cs{vsplit} 产生的盒子的头部插入一个粘连，
% 从而使竖直盒子的顶部不会触发 Underfull 警告。
%    \begin{macrocode}
        \tex_splittopskip:D 0pt plus 1fil
        \vbox_set:Nn \l_@@_tmpa_box
          {
            \tex_vskip:D 0pt plus 1fil
            \tex_hsize:D #1 \centering #3
          }
%    \end{macrocode}
%
% 随后，使用 \cs{vsplit} 命令从先前用于预排版的竖直盒子中竖直切割出每一行，
% 并将其存入一个新的竖直盒子。
%    \begin{macrocode}
        \bool_until_do:nn { \box_if_empty_p:N \l_@@_tmpa_box }
          {
            \vbox_set_split_to_ht:NNn \l_@@_tmpb_box \l_@@_tmpa_box
              { \baselineskip }
%    \end{macrocode}
%
% 对于切割所得到的每个竖直盒子，我们将该盒子的内容重新封装至一个新的给定高度的、
% 基线在盒子顶部的竖直盒子中。需注意的是，为了实现与普通文字相同的效果，我们需
% 在该盒子顶部填充一个垂直间距 \cs{vfil} 以将文字置于盒子底部。另一方面，我们
% 还需要设置盒子的深度（基线到线或点的顶部的距离）。参考 \pkg{xeCJKfntef} 宏包
% 对深度的设置，我们将每个盒子的深度设置为 \qty{0.2}{em}。
%    \begin{macrocode}
            \vbox_set_to_ht:Nnn \l_@@_tmpc_box { #2 }
              { \tex_vfil:D \vbox_unpack:N \l_@@_tmpb_box }
            \box_set_dp:Nn \l_@@_tmpc_box { 0.2 em }
            \box_use:N \l_@@_tmpc_box
%    \end{macrocode}
%
% 最后，在每一行的底部加上一条长度给定、宽度为 \qty{0.5}{pt} 的线段作为下划线。
%    \begin{macrocode}
            \tex_hrule:D width #1 height 0.5pt
          }
      }
  }
%    \end{macrocode}
%
%    \begin{macrocode}
%    \end{macrocode}
% \end{macro}
%
% \begin{macro}{\@@_cover_pg_title:,\@@_cover_pg_title_en:}
% 博士、硕士学位论文封面中外文标题部件。
%    \begin{macrocode}
\cs_new_protected:Npn \@@_cover_pg_title:
  {
    \@@_get_text_width:Nn \l_@@_tmpa_dim 
      { \bfseries \zihao {1} \c_@@_thesis_title_tl \c_@@_colon_tl }
    \group_begin:
      \bfseries \c_@@_thesis_title_tl \c_@@_colon_tl
    \group_end:
    \@@_pg_title_fixed_ht_box:nnn { 9.8 cm } { 1.8 cm }
      { \l_@@_info_title_tl }
  }
\cs_new_protected:Npn \@@_cover_pg_title_en:
  {
    \@@_get_text_width:Nn \l_@@_tmpa_dim 
      {
        \bfseries \zihao {3} \c_@@_thesis_title_en_tl
        \c_@@_colon_en_tl \c_space_token
      }
    \group_begin:
      \c_@@_thesis_title_en_tl \c_@@_colon_en_tl \c_space_token
    \group_end:
    \@@_pg_title_fixed_ht_box:nnn { 12 cm } { 1.2 cm }
      { \l_@@_info_title_en_tl }
  }
%    \end{macrocode}
% \end{macro}
%
% \begin{macro}{\@@_cover_ug_info:}
% 本科生毕业论文封面中的个人信息（姓名、学号、院系、专业、导师及其职称）。
% 在模板中，个人信息的行距被设置为四号字体、 1.5 倍行距。
%    \begin{macrocode}
\cs_new_protected:Npn \@@_cover_ug_info:
  {
    \begin{minipage} [ c ] { \textwidth }
      \centering \zihao { 4 }
      \clist_set:Nx \l_@@_tmpa_clist
        {
          \c_@@_info_name_tl,
          \c_@@_info_student_id_tl,
          \c_@@_info_department_tl,
          \c_@@_info_major_tl,
          \c_@@_info_supervisor_tl,
          \c_@@_info_academic_title_tl
        }
      \clist_set:Nx \l_@@_tmpb_clist
        {
          { \l_@@_info_author_tl        },
          { \l_@@_info_student_id_tl     },
          { \l_@@_info_department_tl    },
          { \l_@@_info_major_tl         },
          { \l_@@_info_supervisor_tl    },
          { \l_@@_info_academic_title_tl }
        }
      \dim_set:Nn \l_@@_tmpb_dim { 6 cm }
      \clist_map_inline:Nn \l_@@_tmpb_clist
        {
          \@@_get_text_width:Nn \l_@@_tmpa_dim { \bfseries ##1 }
          \dim_set:Nn \l_@@_tmpb_dim
            { \dim_max:nn { \l_@@_tmpb_dim } { \l_@@_tmpa_dim } }
        }
      \bool_until_do:nn
        { \clist_if_empty_p:N \l_@@_tmpa_clist }
        {
          \bfseries
          \clist_pop:NN \l_@@_tmpa_clist \l_@@_tmpa_tl
          \clist_pop:NN \l_@@_tmpb_clist \l_@@_tmpb_tl
          \@@_spread_box:no { 4 em } { \l_@@_tmpa_tl }
          \c_@@_colon_tl
          \underline{
            \@@_center_box:nn { \l_@@_tmpb_dim } { \l_@@_tmpb_tl }
          }
          \skip_vertical:n { 0 ex }
        }
    \end{minipage}
  }
%    \end{macrocode}
% \end{macro}
%
% \begin{macro}{\@@_cover_pg_info:,\@@_cover_pg_info_en:}
% 博士、硕士学位论文封面中的个人信息（姓名、学号、院系、专业、导师）。
%    \begin{macrocode}
\cs_new_protected:Npn \@@_cover_pg_info:
  {
    \begin{minipage} [ c ] { \textwidth }
      \centering \zihao { 4 }
      \tl_if_eq:NnTF \l_@@_info_degree_type_tl { professional }
        {
          \clist_set:Nx \l_@@_tmpa_clist
            {
              \c_@@_info_school_tl,
              \c_@@_info_major_professional_tl,
              \c_@@_info_field_professional_tl,
              \c_@@_info_candidate_tl,
              \c_@@_info_supervisor_tl,
            }
        }
        {
          \clist_set:Nx \l_@@_tmpa_clist
            {
              \c_@@_info_school_tl,
              \c_@@_info_major_tl,
              \c_@@_info_field_tl,
              \c_@@_info_candidate_tl,
              \c_@@_info_supervisor_tl,
            }
        }
      \clist_set:Nx \l_@@_tmpb_clist
        {
          { \l_@@_info_department_tl    },
          { \l_@@_info_major_tl         },
          { \l_@@_info_field_tl         },
          { \l_@@_info_author_tl        },
          { \l_@@_info_supervisor_tl    },
        }
      \dim_set:Nn \l_@@_tmpb_dim { 10 em }
      \clist_map_inline:Nn \l_@@_tmpb_clist
        {
          \@@_get_text_width:Nn \l_@@_tmpa_dim { \bfseries ##1 }
          \dim_set:Nn \l_@@_tmpb_dim
            { \dim_max:nn { \l_@@_tmpb_dim } { \l_@@_tmpa_dim } }
        }
      \bool_until_do:nn
        { \clist_if_empty_p:N \l_@@_tmpa_clist }
        {
          \bfseries
          \clist_pop:NN \l_@@_tmpa_clist \l_@@_tmpa_tl
          \clist_pop:NN \l_@@_tmpb_clist \l_@@_tmpb_tl
          \tl_if_eq:NnTF \l_@@_info_degree_type_tl { professional }
            { \@@_spread_box:no { 6 em } { \l_@@_tmpa_tl } }
            { \@@_spread_box:no { 5 em } { \l_@@_tmpa_tl } }
          \c_@@_colon_tl
          \uline{\@@_center_box:nn { \l_@@_tmpb_dim } { \l_@@_tmpb_tl }}
          \skip_vertical:n { 0 ex }
        }
    \end{minipage}
  }
\cs_new_protected:Npn \@@_cover_pg_info_en:
  {
    \begin{minipage} [ c ] { \textwidth }
      \centering \zihao { -4 }
      \tl_if_eq:NnTF \l_@@_info_degree_type_tl { professional }
        {
          \clist_set:Nx \l_@@_tmpa_clist
            {
              \c_@@_info_school_en_tl,
              \c_@@_info_major_pro_en_tl,
              \c_@@_info_field_pro_en_tl,
              \c_@@_info_candidate_en_tl,
              \c_@@_info_supervisor_en_tl,
            }
        }
        {
          \clist_set:Nx \l_@@_tmpa_clist
            {
              \c_@@_info_school_en_tl,
              \c_@@_info_major_en_tl,
              \c_@@_info_field_en_tl,
              \c_@@_info_candidate_en_tl,
              \c_@@_info_supervisor_en_tl,
            }
        }
      \clist_set:Nx \l_@@_tmpb_clist
        {
          { \l_@@_info_department_en_tl    },
          { \l_@@_info_major_en_tl         },
          { \l_@@_info_field_en_tl         },
          { \l_@@_info_author_en_tl        },
          { \l_@@_info_supervisor_en_tl    },
        }
      \dim_set:Nn \l_@@_tmpb_dim { 10 em }
      \dim_set:Nn \l_@@_tmpc_dim { 1 em }
      \clist_map_inline:Nn \l_@@_tmpb_clist
        {
          \@@_get_text_width:Nn \l_@@_tmpa_dim { \bfseries \zihao{-4} ##1 }
          \dim_set:Nn \l_@@_tmpb_dim
            { \dim_max:nn { \l_@@_tmpb_dim } { \l_@@_tmpa_dim } }
        }
      \clist_map_inline:Nn \l_@@_tmpa_clist
        {
          \@@_get_text_width:Nn \l_@@_tmpa_dim
            { \bfseries \zihao{-4} ##1 \c_@@_colon_en_tl \c_space_token }
          \dim_set:Nn \l_@@_tmpc_dim
            { \dim_max:nn { \l_@@_tmpc_dim } { \l_@@_tmpa_dim } }
        }
      \bool_until_do:nn
        { \clist_if_empty_p:N \l_@@_tmpa_clist }
        {
          \bfseries
          \clist_pop:NN \l_@@_tmpa_clist \l_@@_tmpa_tl
          \clist_pop:NN \l_@@_tmpb_clist \l_@@_tmpb_tl
          \mode_leave_vertical:
          \hbox_to_wd:nn { \l_@@_tmpc_dim }
            { \l_@@_tmpa_tl \c_@@_colon_en_tl \c_space_token \tex_hfil:D }
          \uline{\@@_center_box:nn { \l_@@_tmpb_dim } { \l_@@_tmpb_tl }}
          \skip_vertical:n { 0 ex }
        }
    \end{minipage}
  }
%    \end{macrocode}
% \end{macro}
%
% \begin{macro}{\@@_cover_ug_date:,\@@_cover_pg_date:,\@@_cover_pg_date_en:}
% 封面底部的论文完成时间信息。
%    \begin{macrocode}
\cs_new_protected:Npn \@@_cover_ug_date:
  {
    \@@_fixed_width_center_box:nn { \textwidth }
      {
        \bfseries \zihao { 4 } 
        \theyear \c_@@_year_tl \themonth \c_@@_month_tl
      }
  }
\cs_new_protected:Npn \@@_cover_pg_date:
  {
    \@@_fixed_width_center_box:nn { \textwidth }
      {
        \bfseries
        \theyear \c_@@_year_tl \themonth \c_@@_month_tl
      }
  }
\cs_new_protected:Npn \@@_cover_pg_date_en:
  {
    \@@_fixed_width_center_box:nn { \textwidth }
      {
        \bfseries
        \clist_item:Nn \c_@@_month_en_clist { \themonth }
        , \c_space_token \theyear
      }
  }
%    \end{macrocode}
% \end{macro}
%
% \begin{macro}{\@@_orig_text_pg:}
% 原创性声明内容。
%    \begin{macrocode}
\cs_new_protected:Npn \@@_orig_text_pg:
  {
    \indent \c_@@_orig_text_pg_part_one_tl
    \c_@@_left_guillemet_tl
    \l_@@_info_title_trim_tl
    \c_@@_right_guillemet_tl
    \c_@@_orig_text_pg_part_two_tl
  }
%    \end{macrocode}
% \end{macro}
%
% \begin{macro}{\@@_auth_text_pg:}
% 论文使用授权内容。
%    \begin{macrocode}
\cs_new_protected:Npn \@@_auth_text_pg:
  {
    \indent
    \c_@@_left_guillemet_tl \l_@@_info_title_trim_tl
    \c_@@_right_guillemet_tl
    \c_@@_auth_text_pg_part_one_tl \par
    \c_@@_auth_text_pg_part_two_tl \par
    \c_@@_left_round_bracket_tl
    \hbox_to_wd:nn { 2.5 em } { } \c_@@_right_round_bracket_tl
    \c_@@_auth_text_pg_option_one_tl \par
    \c_@@_left_round_bracket_tl
    \hbox_to_wd:nn { 2.5 em } { } \c_@@_right_round_bracket_tl
    \c_@@_auth_text_pg_option_two_tl \par
  }
%    \end{macrocode}
% \end{macro}
%
% \begin{macro}{\@@_orig_sign_area_ug:,\@@_orig_sign_area_pg:}
% 原创性声明签名。
%    \begin{macrocode}
\cs_new_protected:Npn \@@_orig_sign_area_ug:
  {
    \c_@@_orig_sign_tl \c_@@_colon_tl
    \hbox_to_wd:nn { 15 em } { }
    \c_@@_date_tl \c_@@_colon_tl
    \hbox_to_wd:nn { 3.4 em } { } \c_@@_year_tl
    \hbox_to_wd:nn { 1.5 em } { } \c_@@_month_tl
    \hbox_to_wd:nn { 1.5 em } { } \c_@@_day_tl
  }
\cs_new_protected:Npn \@@_orig_sign_area_pg:
  {
    \indent
    { \bfseries \c_@@_author_sign_tl } \hbox_to_wd:nn { 15.5 em } { }
    { \bfseries \c_@@_date_tl }
    \hbox_to_wd:nn { 4 em } { } \c_@@_year_tl
    \hbox_to_wd:nn { 2 em } { } \c_@@_month_tl
    \hbox_to_wd:nn { 2 em } { } \c_@@_day_tl
  }
%    \end{macrocode}
% \end{macro}
%
% \begin{macro}{\@@_auth_sign_area_ug:,\@@_auth_sign_area_pg:}
% 论文使用授权声明。
%    \begin{macrocode}
\cs_new_protected:Npn \@@_auth_sign_area_ug:
  {
    \c_@@_author_sign_tl \c_@@_colon_tl
    \hbox_to_wd:nn { 5.5 em } { }
    \c_@@_supervisor_sign_tl \c_@@_colon_tl
    \hbox_to_wd:nn { 5.5 em } { }
    \c_@@_date_tl \c_@@_colon_tl
    \hbox_to_wd:nn { 3.4 em } { } \c_@@_year_tl
    \hbox_to_wd:nn { 1.5 em } { } \c_@@_month_tl
    \hbox_to_wd:nn { 1.5 em } { } \c_@@_day_tl
  }
\cs_new_protected:Npn \@@_auth_sign_area_pg:
  {
    \indent
    { \bfseries \c_@@_supervisor_sign_tl } \hfill
    { \bfseries \c_@@_author_sign_tl } \hbox_to_wd:nn { 14 em } { } \par
    \hfill
    { \bfseries \c_@@_date_tl }
    \hbox_to_wd:nn { 4 em } { } \c_@@_year_tl
    \hbox_to_wd:nn { 2 em } { } \c_@@_month_tl
    \hbox_to_wd:nn { 2 em } { } \c_@@_day_tl
    \hbox_to_wd:nn { 5 em } { }
  }
%    \end{macrocode}
% \end{macro}
%
% \begin{macro}{style/logo-src}
% 定义封面 logo 接口。
%    \begin{macrocode}
\tl_new:N \l_@@_logo_resource_tl
\keys_define:nn { ecnu / style }
  {
    logo-src  .tl_set:N   = \l_@@_logo_resource_tl,
    logo-src  .initial:n  =
      {
        \bool_if:NTF \g_@@_is_bachelor_bool
          { ecnu-vi-inner-cover-logo-ug.pdf }
          { ecnu-vi-inner-cover-logo-pg.pdf }
      }
  }
%    \end{macrocode}
% \end{macro}
%
% \begin{macro}{\@@_cover_pg_logo_title:,\@@_cover_pg_logo_title_en:}
% 博士、硕士学位论文中封面 logo 部分的中外文标题。
%    \begin{macrocode}
\cs_new_protected:Npn \@@_cover_pg_logo_title:
  { 
    \prop_get:NVN \c_@@_thesis_type_prop
      \g_@@_thesis_type_tl \l_@@_tmpa_tl
    \l_@@_tmpa_tl
    \tl_if_eq:NnT \l_@@_info_degree_type_tl { professional }
      { \c_@@_professional_tl }
    \c_@@_degree_tl \c_@@_thesis_tl
  }
\cs_new_protected:Npn \@@_cover_pg_logo_title_en:
  { 
    \tl_if_eq:NnTF \l_@@_info_degree_type_tl { professional }
      {
        \tl_if_eq:NnTF \g_@@_thesis_type_tl { master }
          {
            \c_@@_masters_en_tl \c_space_token \c_@@_degree_en_tl
            \c_space_token \c_@@_thesis_en_tl
            \c_space_token (\c_@@_professional_en_tl)
          }
          {
            \text_uppercase:n { \c_@@_doctoral_en_tl }
            \c_space_token
            \text_uppercase:n { \c_@@_dissertation_en_tl }
          }
      }
      {
        \tl_if_eq:NnTF \g_@@_thesis_type_tl { master }
          { \text_uppercase:n { \c_@@_masters_en_tl } }
          { \text_uppercase:n { \c_@@_doctoral_en_tl } }
        \c_space_token
        \text_uppercase:n { \c_@@_dissertation_en_tl }
      }
  }
%    \end{macrocode}
% \end{macro}
%
% \begin{macro}{\@@_cover_ug_logo:,\@@_cover_pg_logo:,\@@_cover_pg_logo_en:}
% 封面 logo 部件。
%    \begin{macrocode}
\cs_new_protected:Npn \@@_cover_ug_logo:
  { 
    \tl_if_empty:NF \l_@@_logo_resource_tl
      {
        \begin{minipage} [ c ] { \textwidth }
          \centering
          \includegraphics{\l_@@_logo_resource_tl}
        \end{minipage}
      }
  }
\cs_new_protected:Npn \@@_cover_pg_logo:
  { 
    \begin{minipage} [ c ] { \textwidth }
      \@@_line_spread:N \c_@@_single_space_line_spread_fp
      \centering \zihao { 5 } 
      \tl_if_empty:NF \l_@@_logo_resource_tl
        {
          \includegraphics[width=8cm]{\l_@@_logo_resource_tl}
          \skip_vertical:n { 0 em }
        }
      \c_@@_university_name_en_tl
      \skip_vertical:n { 5 bp }
      \group_begin:
        \bfseries
        \@@_cover_pg_logo_title:
        \skip_vertical:n { 1.12 bp }
        \@@_cover_pg_logo_title_en:
      \group_end:
    \end{minipage}
  }
\cs_new_protected:Npn \@@_cover_pg_logo_en:
  { 
    \begin{minipage} [ c ] { \textwidth }
      \@@_line_spread:N \c_@@_single_space_line_spread_fp
      \centering \zihao { 5 } \bfseries
      \c_@@_university_name_en_tl
    \end{minipage}
  }
%    \end{macrocode}
% \end{macro}
%
% \subsubsection{封面实现接口}
%
%    \begin{macrocode}
\tl_new:N \l_@@_cover_template_tl
\DeclareObjectType { ecnu / cover } { \c_zero_int }
\cs_new_protected:Npn \@@_cover_declare_template_interface:nn #1#2
  { \DeclareTemplateInterface { ecnu / cover } {#1} { \c_zero_int } {#2} }
\cs_new_protected:Npn \@@_cover_declare_template_code:nnn #1#2#3
  { \DeclareTemplateCode { ecnu / cover } {#1} { \c_zero_int } {#2} {#3} }
\cs_generate_variant:Nn \@@_cover_declare_template_interface:nn { nx  }
\cs_generate_variant:Nn \@@_cover_declare_template_code:nnn     { nxn }
\cs_new:Npn \@@_cover_key_type:n #1
  {
    #1 / content     : tokenlist,
    #1 / format      : tokenlist,
    #1 / bottom-skip : skip,
    #1 / align       : choice { left, right, center, normal } = normal,
  }
\cs_new:Npn \@@_cover_key_binding:n #1
  {
    #1 / content     =
      \exp_not:c
        { l_@@_cover / \l_@@_cover_template_tl / #1 / content_tl  },
    #1 / format      =
      \exp_not:c
        { l_@@_cover / \l_@@_cover_template_tl / #1 / format_tl   },
    #1 / bottom-skip =
      \exp_not:c
        { l_@@_cover / \l_@@_cover_template_tl / #1 / bottom_skip },
    #1 / align       =
      {
        left   =
          \exp_not:N \cs_set_protected:cpn
            { @@_cover / \l_@@_cover_template_tl / #1 / align:n }
            \exp_not:n {##1}
            {
              \exp_not:n
                {
                  \group_begin:
                    \raggedright ##1
                  \group_end:
                }
            },
        right  =
          \exp_not:N \cs_set_protected:cpn
            { @@_cover / \l_@@_cover_template_tl / #1 / align:n }
            \exp_not:n {##1}
            {
              \exp_not:n
                {
                  \group_begin:
                    \raggedleft ##1
                  \group_end:
                }
            },
        center =
          \exp_not:N \cs_set_protected:cpn
            { @@_cover / \l_@@_cover_template_tl / #1 / align:n }
            \exp_not:n {##1}
            {
              \exp_not:n
                {
                  \group_begin:
                    \centering ##1
                  \group_end:
                }
            },
        normal =
          \exp_not:N \cs_set_protected:cpn
            { @@_cover / \l_@@_cover_template_tl / #1 / align:n }
            \exp_not:n {##1}
            { \exp_not:n { \group_begin: ##1 \group_end: } }
      },
  }
\cs_new_protected:Npn \@@_cover_declare_template:nn #1#2
  {
    \tl_set:Nn \l_@@_cover_template_tl {#1}
    \@@_cover_declare_template_interface:nx {#1}
      {
        format      : tokenlist,
        top-skip    : skip,
        bottom-skip : skip,
        \clist_map_function:nN {#2} \@@_cover_key_type:n
      }
    \tl_new:c   { l_@@_cover / #1 / format_tl   }
    \skip_new:c { l_@@_cover / #1 / top_skip    }
    \skip_new:c { l_@@_cover / #1 / bottom_skip }
    \clist_map_inline:nn {#2}
      {
        \tl_new:c   { l_@@_cover / #1 / ##1 / content_tl  }
        \tl_new:c   { l_@@_cover / #1 / ##1 / format_tl   }
        \skip_new:c { l_@@_cover / #1 / ##1 / bottom_skip }
      }
    \@@_cover_declare_template_code:nxn {#1}
      {
        format      = \exp_not:c { l_@@_cover / #1 / format_tl   },
        top-skip    = \use:c     { l_@@_cover / #1 / top_skip    },
        bottom-skip = \use:c     { l_@@_cover / #1 / bottom_skip },
        \clist_map_function:nN {#2} \@@_cover_key_binding:n
      }
      {
        \AssignTemplateKeys
        \tl_use:c        { l_@@_cover / #1 / format_tl }
        \clist_map_inline:nn {#2}
          { 
            \noindent
            \use:c { @@_cover / #1 / ####1 / align:n }
              {
                \tl_use:c { l_@@_cover / #1 / ####1 / format_tl  }
                \tl_use:c { l_@@_cover / #1 / ####1 / content_tl }
                \par
              }
            \@@_vspace:c { l_@@_cover / #1 / ####1 / bottom_skip }
          }
        \@@_vspace:c { l_@@_cover / #1 / bottom_skip }
        \clearpage{\thispagestyle{empty}\cleardoublepage}
      }
  }
\NewDocumentCommand \DeclareCoverTemplate { m m }
  { \@@_cover_declare_template:nn {#1} {#2} }
\DeclareCoverTemplate { cover-i   }
  { top, logo, title, title-en, info, date }
\DeclareCoverTemplate { cover-ii  }
  {
    originality-title,
    originality-text,
    originality-sign,
    authorization-title,
    authorization-text,
    authorization-sign,
    others,
  }
\DeclareInstance { ecnu / cover } { cover-i-ug } { cover-i }
  {
    bottom-skip            = 0 pt plus 0.5 fill,
    top      / content     = \@@_cover_ug_top:,
    logo     / content     = \@@_cover_ug_logo:,
    title    / content     =
      \expandafter\uline\expandafter{\l_@@_info_title_tl},
    title-en / content     =
      \expandafter\uline\expandafter{\l_@@_info_title_en_tl},
    info     / content     = \@@_cover_ug_info:,
    date     / content     = \@@_cover_ug_date:,
    title    / format      =
      \@@_line_spread:n { 1.1 } \zihao { 1 } \bfseries,
    title-en / format      =
      \@@_line_spread:n { 1.1 } \zihao { 1 } \bfseries,
    top      / bottom-skip = 0 pt plus 0.5 fill,
    logo     / bottom-skip = 0 pt plus 0.3 fill,
    title-en / bottom-skip = 0 pt plus 1 fill,
    info     / bottom-skip = 31.2 bp,
    top      / align       = right,
    logo     / align       = center,
    title    / align       = center,
    title-en / align       = center,
    info     / align       = center,
    date     / align       = center,
  }
\DeclareInstance { ecnu / cover } { cover-i-pg } { cover-i }
  {
    bottom-skip            = 0 pt,
    top      / content     = \@@_cover_pg_top:,
    logo     / content     = \@@_cover_pg_logo:,
    title    / content     = \@@_cover_pg_title:,
    info     / content     = \@@_cover_pg_info:,
    date     / content     = \@@_cover_pg_date:,
    title    / format      = \zihao { 1 } \bfseries,
    top      / format      = \zihao { 5 },
    top      / bottom-skip = 25.272 bp,
    logo     / bottom-skip = 1 em plus .5 fill,
    title    / bottom-skip = 1 em plus 1 fill,
    info     / bottom-skip = 2 em plus .5 fill,
    top      / align       = center,
    logo     / align       = center,
    title    / align       = center,
    info     / align       = center,
    date     / align       = center,
  }
  \DeclareInstance { ecnu / cover } { cover-i-pg-en } { cover-i }
  {
    bottom-skip            = 0 pt,
    top      / content     = \@@_cover_pg_top_en:,
    logo     / content     = \@@_cover_pg_logo_en:,
    title    / content     = \@@_cover_pg_title_en:,
    info     / content     = \@@_cover_pg_info_en:,
    date     / content     = \@@_cover_pg_date_en:,
    title    / format      = \zihao { 3 } \bfseries,
    top      / format      = \zihao { 5 },
    top      / bottom-skip = 49.696 bp,
    logo     / bottom-skip = 1 em plus .5 fill,
    title    / bottom-skip = 1 em plus 1 fill,
    info     / bottom-skip = 2 em plus .5 fill,
    top      / align       = center,
    logo     / align       = center,
    title    / align       = center,
    info     / align       = center,
    date     / align       = center,
  }
\DeclareInstance { ecnu / cover } { cover-ii-ug } { cover-ii }
  {
    originality-title   / content     = \c_@@_orig_title_ug_tl,
    originality-text    / content     = \c_@@_orig_text_ug_tl,
    originality-sign    / content     = \@@_orig_sign_area_ug:,
    authorization-title / content     = \c_@@_auth_title_ug_tl,
    authorization-text  / content     = \c_@@_auth_text_ug_tl,
    authorization-sign  / content     = \@@_auth_sign_area_ug:,
    originality-title   / format      =
      \@@_line_spread:n { 1.85714286 } \zihao { 4 } \bfseries,
    originality-text    / format      = \@@_line_spread:n { 1.625 },
    authorization-title / format      =
      \@@_line_spread:n { 1.85714286 } \zihao { 4 } \bfseries,
    authorization-text  / format      = \@@_line_spread:n { 1.625 },
    originality-sign    / bottom-skip = 70.2 bp,
    authorization-text  / bottom-skip = 23.4 bp,
    originality-title   / align       = center,
    originality-sign    / align       = right,
    authorization-title / align       = center,
    authorization-sign  / align       = right,
  }
\DeclareInstance { ecnu / cover } { cover-ii-pg } { cover-ii }
  {
    format                            = 
      \@@_line_spread:n { 1.85714286 } \zihao { 5 },
    originality-title   / content     = \c_@@_orig_title_pg_tl,
    originality-text    / content     = \@@_orig_text_pg:,
    originality-sign    / content     = \@@_orig_sign_area_pg:,
    authorization-title / content     = \c_@@_auth_title_pg_tl,
    authorization-text  / content     = \@@_auth_text_pg:,
    authorization-sign  / content     = \@@_auth_sign_area_pg:,
    others              / content     = 
      \indent \c_@@_auth_others_text_pg_tl,
    originality-title   / format      = \zihao { 4 } \bfseries,
    originality-text    / format      = \zihao { 5 },
    others              / format      = 
      \@@_line_spread:n { 2.1388889 } \zihao { -5 },
    authorization-title / format      = \zihao { 4 } \bfseries,
    authorization-text  / format      = \zihao { 5 },
    originality-title   / bottom-skip = 6 bp,
    originality-text    / bottom-skip = 3 bp,
    originality-sign    / bottom-skip = 55.8 bp,
    authorization-title / bottom-skip = 6 bp,
    authorization-text  / bottom-skip = 3 bp,
    authorization-sign  / bottom-skip = 29.4 bp,
    originality-title   / align       = center,
    authorization-title / align       = center,
  }
\tl_new:N \l_@@_declaration_page_tl
\keys_define:nn { ecnu / style }
  {
    declaration-src   .code:n =
      {
        \tl_set_eq:NN \l_@@_declaration_page_tl \l_keys_value_tl
        \RequirePackage { pdfpages }
      },
  }
%    \end{macrocode}
%
% \begin{macro}{\innercover,\innercoverEN}
% 中外文内封面命令。
%    \begin{macrocode}
\NewDocumentCommand \innercover { }
  {
    \thispagestyle { empty }
    \tl_set:Nn \thepage { A }
    \bool_if:NTF \g_@@_is_bachelor_bool
      { \UseInstance { ecnu / cover } { cover-i-ug } }
      { \UseInstance { ecnu / cover } { cover-i-pg } }
  }
\NewDocumentCommand \innercoverEN { }
  {
    \thispagestyle { empty }
    \tl_set:Nn \thepage { A-EN }
    \UseInstance { ecnu / cover } { cover-i-pg-en }
  }
%    \end{macrocode}
% \end{macro}
%    
% \begin{macro}{\declarationpage}
% 声明页命令。
%    \begin{macrocode}
\NewDocumentCommand \declarationpage { }
  {
    \cleardoublepage
    \thispagestyle { empty }
    \tl_set:Nn \thepage { B }
    \tl_if_empty:NTF \l_@@_declaration_page_tl
      {
        \bool_if:NTF \g_@@_is_bachelor_bool
          { \UseInstance { ecnu / cover } { cover-ii-ug } }
          { \UseInstance { ecnu / cover } { cover-ii-pg } }
      }
      { \includepdf [pages=-] { \l_@@_declaration_page_tl } }
  }
%    \end{macrocode}
% \end{macro}
%
% \subsection{目录}
%
% 设置目录标题以及目录层级。目录层级仅显示一级和二级标题。
%    \begin{macrocode}
\keys_set:nn { ctex }
  {
    contentsname    = \c_@@_name_toc_tl,
    listfigurename  = \c_@@_name_lof_tl,
    listtablename   = \c_@@_name_lot_tl,
    tocdepth        = \bool_if:NTF \g_@@_is_bachelor_bool { 1 } { 2 },
  }
%    \end{macrocode}
%
% 使用 \pkg{titletoc} 设置目录格式。
%    \begin{macrocode}
\RequirePackage { titletoc }
%    \end{macrocode}
% 
% \begin{macro}{\tocchapterformat,\tocsectionformat}
% 单独定义命令配置目录中章节标题、节标题的具体内容。其中，\cs{contentsmargin} 
% 用于设置右边距。我们将其设置为 \qty{0}{pt}，以去除点线与页码之间的间距。
% 
% 在本科生毕业论文中，章标题的字体为宋体加粗，字号为小四；节标题的字体
% 为宋体，字号为小四。章标题上方有 \qty{6}{pt} 的垂直间隙。
% 
% 在博士、硕士学位论文中，章标题的字体为黑体，字号为四号；节标题的字体
% 为黑体，字号为小四。
%    \begin{macrocode}
\NewDocumentCommand \tocchapterformat { }
  {
    \addvspace{6pt} \contentsmargin{0pt}
    \bool_if:NTF \g_@@_is_bachelor_bool
      { \rmfamily \bfseries }
      { \sffamily \zihao{ 4 } }
  }
\NewDocumentCommand \tocsectionformat {}
  {
    \contentsmargin{0pt}
    \bool_if:NF \g_@@_is_bachelor_bool { \sffamily }
  }
%    \end{macrocode}
% \end{macro}
%
% \begin{macro}{\tocchaptercontents,\tocsectioncontents,\tocsubsectioncontents}
% 设置目录数字编号格式。
%    \begin{macrocode}
\NewDocumentCommand \tocchaptercontents { }
  {
    \thecontentslabel
    \bool_if:NF \g_@@_is_bachelor_bool { \hspace*{2ex} }
  }
\NewDocumentCommand \tocsectioncontents { }
  {
    \thecontentslabel
    \bool_if:NTF \g_@@_is_bachelor_bool
      { \tl_if_eq:NnT \g_@@_numbering_tl { arabic } { \hspace*{1ex} } }
      { \hspace*{2ex} }
  }
\NewDocumentCommand \tocsubsectioncontents { }
  {
    \thecontentslabel
    \tl_if_eq:NnT \g_@@_numbering_tl { arabic } { \hspace*{2ex} }
  }
%    \end{macrocode}
% \end{macro}
%
% 调整目录各层级项目格式。
%    \begin{macrocode}
\titlecontents{chapter}[0em]
              {\tocchapterformat}{\tocchaptercontents}{}
              {\titlerule*[.25pc]{.}{\rmfamily\thecontentspage}}
              [\addvspace{3pt}]
\titlecontents{section}[1em]
              {\tocsectionformat}{\tocsectioncontents}{}
              {\titlerule*[.25pc]{.}{\rmfamily\thecontentspage}}
              [\addvspace{3pt}]
\titlecontents{subsection}[2em]
              {\contentsmargin{0pt}}{\tocsubsectioncontents}{}
              {\titlerule*[.25pc]{.}{\rmfamily\thecontentspage}}
%    \end{macrocode}
%
% 调整插图目录和随表目录的项目格式。
%    \begin{macrocode}
\titlecontents{figure}[0em]{\contentsmargin{0pt}}
              {\figurename\thecontentslabel\hspace*{2ex}}{}
              {\titlerule*[.25pc]{.}\thecontentspage}
\titlecontents{table}[0em]{\contentsmargin{0pt}}
              {\tablename\thecontentslabel\hspace*{2ex}}{}
              {\titlerule*[.25pc]{.}\thecontentspage}
%    \end{macrocode}
%
% \begin{macro}{\@@_toc_page:Nn}
% 自定义标题页格式。来自于 \LaTeXe{} 标准文档类 \file{book.cls}。
%    \begin{macrocode}
\cs_new_protected:Npn \@@_toc_page:Nn #1#2
  {
    \group_begin:
%    \end{macrocode}
%
% 在本科生毕业论文中，目录标题的大小为小四、字体为宋体加粗。
%    \begin{macrocode}
      \bool_if:NT \g_@@_is_bachelor_bool
        {
          \pagestyle{ empty }
          \ctexset
            {
              chapter / pagestyle   = empty,
              chapter / numbering   = false,
              chapter / format      = \rmfamily\bfseries\centering,
            }
        }
%    \end{macrocode}
% 
%    \begin{macrocode}
      \chapter* { #1 }
      \bool_if:NTF \g_@@_is_bachelor_bool
        { \@@_line_spread:N \c_@@_single_space_line_spread_fp }
        { \@@_line_spread:n { 1.3 } }
      \@starttoc{#2}
    \group_end:
    \cleardoublepage
    \pagestyle{ fancy }
  }
%    \end{macrocode}
% \end{macro}
% 
% \begin{macro}{\tableofcontents,\listoffigures,\listoftables}
% 修改 \cs{tableofcontents}、\cs{listoffigures}、\cs{listoftables}。
%    \begin{macrocode}
\RenewDocumentCommand \tableofcontents { }
  { \@@_toc_page:Nn \contentsname {toc} }
\RenewDocumentCommand \listoffigures { }
  { \@@_toc_page:Nn \listfigurename {lof} }
\RenewDocumentCommand \listoftables { }
    { \@@_toc_page:Nn \listtablename {lot} }
%    \end{macrocode}
% \end{macro}
%
% \subsection{摘要}
%
% \begin{environment}{abstract}
% \begin{environment}{abstractEN}
% 定义中外文摘要环境。
%    \begin{macrocode}
\RenewDocumentEnvironment { abstract  } { }
  { \@@_abstract_begin: } { \@@_abstract_end: }
\NewDocumentEnvironment { abstractEN } { }
  { \@@_abstract_en_begin: } { \@@_abstract_en_end: }
%    \end{macrocode}
% \end{environment}
% \end{environment}
%
% \begin{macro}{\@@_abstract_begin:,\@@_abstract_en_begin:}
% 中外文摘要页环境起始部分。在环境起始部分手动添加中外文标题。
%
% 在本科生毕业论文中，需将“摘要”（“Abstract”）添加至目录；摘要部分字体大小为五号。
%    \begin{macrocode}
\cs_new_protected:Npn \@@_abstract_begin:
  {
    \bool_if:NTF \g_@@_is_bachelor_bool
      {
        \@@_abs_title:V \l_@@_info_title_trim_tl
        \addcontentsline{toc}{chapter}{\c_@@_name_abstract_tl}
        \small
        \group_begin:
          \skip_vertical:n { 17.6 bp }
          \noindent \rmfamily \bfseries
          \c_@@_name_abstract_tl \c_@@_colon_tl
          \skip_vertical:n { 3 bp }
        \group_end:
        \par
      }
      { \@@_abs_title:V \c_@@_name_chinese_abstract_tl }
  }
\cs_new_protected:Npn \@@_abstract_en_begin:
  {
    \bool_if:NTF \g_@@_is_bachelor_bool
      {
        \@@_absEN_title:V \l_@@_info_title_en_trim_tl
        \addcontentsline{toc}{chapter}{ABSTRACT}
        \small
        \group_begin:
          \skip_vertical:n { 17.6 bp }
          \noindent \rmfamily \bfseries
          \c_@@_name_abstract_en_tl \c_@@_colon_en_tl
          \skip_vertical:n { 3 bp }
        \group_end:
        \par
      }
      {
        \@@_absEN_title:n
          { \text_uppercase:n { \c_@@_name_abstract_en_tl } }
      }
  }
%    \end{macrocode}
% \end{macro}
%
% \begin{macro}{\@@_abstract_end:,\@@_abstract_en_end:}
% 摘要页标题环境结束部分。在摘要部分正文末尾附上关键词。
%    \begin{macrocode}
\cs_new_protected:Npn \@@_abstract_end:
  {
    \@@_keywords:nNnn
      {
        \bool_if:NTF \g_@@_is_bachelor_bool
          { \rmfamily\bfseries }
          { \zihao{4}\sffamily\bfseries }
        \c_@@_name_keywords_tl \c_@@_colon_tl
      }
      \l_@@_info_keywords_clist { ， } { }
    \cleardoublepage
  }
\cs_new_protected:Npn \@@_abstract_en_end:
  {
    \@@_keywords:nNnn
      {
        \bool_if:NTF \g_@@_is_bachelor_bool
          { \rmfamily\bfseries } { \zihao{4}\bfseries }
        \c_@@_name_keywords_en_tl \c_@@_colon_en_tl \c_space_token
      }
      \l_@@_info_keywords_en_clist { ,\c_space_token }
      { \bool_if:NF \g_@@_is_bachelor_bool { \itshape } }
    \cleardoublepage
  }
%    \end{macrocode}
% \end{macro}
%
% \begin{macro}{\@@_keywords:nNnn}
% 中外文摘要页关键词。
%    \begin{macrocode}
\cs_new_protected:Npn \@@_keywords:nNnn #1#2#3#4
  {
    \par \@@_vspace:n { \skip_use:N \baselineskip minus 1 em }
    \par \noindent
    \group_begin: #1 \group_end:
    \group_begin:
      #4
      \clist_use:Nn #2 {#3}
    \group_end:
  }
%    \end{macrocode}
% \end{macro}
%
%
% \subsection{附录}
%
% \begin{environment}{appendix}
% 定义附录环境。
%    \begin{macrocode}
\RenewDocumentEnvironment { appendix } { }
  { \@@_appendix_begin: } { \@@_appendix_end: }
%    \end{macrocode}
% \end{environment}
%
% \begin{macro}{\@@_appendix_begin:}
% 添加附录标题，将附录环境内正文的字号调整为五号。
%    \begin{macrocode}
\cs_new_protected:Npn \@@_appendix_begin:
  { 
    \cleardoublepage
    \@@_title:V \c_@@_name_appendix_tl
    \setcounter{section}{0}
    \bool_if:NT \g_@@_is_bachelor_bool { \small }
  }
%    \end{macrocode}
% \end{macro}
%
% \begin{macro}{\@@_appendix_end:}
% 在打印模式下，保证附录部分为偶数页。
%    \begin{macrocode}
\cs_new_protected:Npn \@@_appendix_end:
  { \cleardoublepage }
%    \end{macrocode}
% \end{macro}
%
%
% \subsection{致谢}
%
% \begin{environment}{acknowledgement}
% 致谢环境。
%    \begin{macrocode}
\NewDocumentEnvironment { acknowledgement } { }
  { \@@_acknowledgement_begin: }
  { \@@_acknowledgement_end: }
%    \end{macrocode}
% \end{environment}
%
% \begin{macro}{\@@_acknowledgement_begin:}
% 致谢页环境。
%    \begin{macrocode}
\cs_new_protected:Npn \@@_acknowledgement_begin:
  {
    \cleardoublepage
    \bool_if:NTF \g_@@_is_bachelor_bool
      { \@@_title:V \c_@@_name_acknowledgement_tl }
      { \@@_title:V \c_@@_name_epilogue_tl }
  }
%    \end{macrocode}
% \end{macro}
%
% \begin{macro}{\@@_acknowledgement_end:}
% 在打印模式下，保证致谢页为偶数页。
%    \begin{macrocode}
\cs_new_protected:Npn \@@_acknowledgement_end:
  { \cleardoublepage }
%    \end{macrocode}
% \end{macro}
%
% \subsection{参考文献}
%
% 模版使用 \cs{biblatex} 和 Biber 处理参考文献。
%
% \begin{variable}{\l_@@_bib_resource_tl}
% 声明变量。
%    \begin{macrocode}
\tl_new:N \l_@@_bib_resource_tl
%    \end{macrocode}
% \end{variable}
%
% \begin{macro}{style/bib-src}
% 定义参考文献数据库路径。
%    \begin{macrocode}
\keys_define:nn { ecnu / style }
  {
    bib-src .tl_set:N = \l_@@_bib_resource_tl
  }
%    \end{macrocode}
% \end{macro}
%
% 由于 \pkg{biblatex} 会写入 begindocument/before 钩子，
% 因此需要在其之前通过写入 env/document/begin
% 钩子载入 \pkg{biblatex} 宏包。
%    \begin{macrocode}
\@@_gadd_ltxhook:nn { env/document/begin }
{
  \@@_biblatex_pre_setup:
  \RequirePackage { biblatex }
  \@@_biblatex_post_setup:
}
%    \end{macrocode}
%
% \begin{macro}{\@@_biblatex_pre_setup:, \@@_biblatex_post_setup:}
% \pkg{biblatex} 宏包相关设置。将参考文献处理后端设置为 biber，
% 参考文献样式遵循 GB/T 7714-2015 标准。
%    \begin{macrocode}
\cs_new_protected:Npn \@@_biblatex_pre_setup:
  {
    \PassOptionsToPackage 
      {
        backend       = biber,
        style         = gb7714-2015,
        seconds       = true,
        hyperref      = true,
        gbpub         = false,
        gbpunctin     = true,
      } 
      { biblatex }
  }
%    \end{macrocode}
% 
% 设置参考文献相关格式。
%    \begin{macrocode}
\cs_new_protected:Npn \@@_biblatex_post_setup:
  {
    \exp_args:NV \addbibresource \l_@@_bib_resource_tl
    \@@_biblatex_allow_url_break:
%    \end{macrocode}
%
% 在本科生毕业论文中，参考文献的字体大小为五号。
%    \begin{macrocode}
    \bool_if:NT \g_@@_is_bachelor_bool
      { \RenewDocumentCommand \bibfont { } { \small } }
%    \end{macrocode}
%
% 参考文献间没有额外间距。
%    \begin{macrocode}
    \setlength{\bibitemsep}{0pt}
    \defbibheading { bibliography } [ \bibname ] 
      { \@@_title:n {##1} }
  }
%    \end{macrocode}
% \end{macro}
%
% \begin{macro}{\@@_biblatex_allow_url_break}
% \pkg{biblatex} 下允许 URL 在字母、数字和一些特殊符号处断行。
%    \begin{macrocode}
\cs_new:Npn \@@_biblatex_allow_url_break:
  {
    \int_set_eq:NN \c@biburlucpenalty  \c_one_int
    \int_set_eq:NN \c@biburlnumpenalty \c_one_int
    \int_set_eq:NN \c@biburllcpenalty  \c_one_int
  }
%    \end{macrocode}
% \end{macro}
%
% \subsection{\pkg{hyperref} 配置}
%
% 在导言区末尾引入 \pkg{hyperref} 宏包，以尽可能减少与其他宏包
% 间的冲突。
%    \begin{macrocode}
\ctex_at_end_preamble:n
  {
    \RequirePackage { hyperref }
    \hypersetup
      {
        hidelinks,
        unicode         = true,
%    \end{macrocode}
% 
% 填写 PDF 元信息。
%    \begin{macrocode}
        pdftitle        = \l_@@_info_title_trim_tl,
        pdfauthor       = \l_@@_info_author_tl,
        pdfkeywords     = \l_@@_info_keywords_clist,
        pdfcreator      = \c_@@_name_pdf_creator_tl
      }
  }
%    \end{macrocode}
%
% \subsection{前导页}
%
% 当设置开启生成封面页或声明页时，模板会在文档前导部分插入前导页（包括封面页和声明页）。
%    \begin{macrocode}
\bool_lazy_or:nnT
  { \g_@@_cover_page_bool } { \g_@@_declaration_page_bool }
  {
    \AtBeginDocument
      {
%    \end{macrocode}
%
% 在本科生毕业论文中，前导页的页边距与正文有不同。
%    \begin{macrocode}
        \bool_if:NT \g_@@_is_bachelor_bool
          {
            \newgeometry
              {
                top     = 2.0 cm, 
                bottom  = 2.0 cm,
                left    = 3.0 cm, 
                right   = 3.0 cm
              }
          }
%    \end{macrocode}
%
% 设置前导页内容。
%    \begin{macrocode}
        \begin{titlepage}
          \bool_if:NT \g_@@_cover_page_bool
            {
              \innercover
              \bool_if:NF \g_@@_is_bachelor_bool { \newpage \innercoverEN }
            }
          \bool_if:NT \g_@@_declaration_page_bool
            { \declarationpage }
        \end{titlepage}
%    \end{macrocode}
%
% 在本科生毕业论文中，前导页设置完成后需要将页面恢复到初始设置。
%    \begin{macrocode}
        \bool_if:NT \g_@@_is_bachelor_bool { \restoregeometry }
      }
  }
%    \end{macrocode}
%
% \subsection{用户接口}
%
% \begin{macro}{info,style}
% 定义元（meta）键值对。
%    \begin{macrocode}
\keys_define:nn { ecnu }
  {
    info  .meta:nn = { ecnu / info  } {#1},
    style .meta:nn = { ecnu / style } {#1}
  }
%    \end{macrocode}
% \end{macro}
%
% \begin{macro}{\ecnusetup}
% 用户设置接口。
%    \begin{macrocode}
\NewDocumentCommand \ecnusetup { m }
  { \keys_set:nn { ecnu } {#1} }
%    \end{macrocode}
% \end{macro}
%
% \begin{macro}{\frontmatter}
% 前置部分开始命令。
%    \begin{macrocode}
\NewDocumentCommand \frontmatter { }
  { \cleardoublepage\pagenumbering{Roman} }
%    \end{macrocode}
% \end{macro}
%
% \begin{macro}{\mainmatter}
% 论文主体部分开始命令。
%    \begin{macrocode}
\NewDocumentCommand \mainmatter { }
  { \cleardoublepage\pagenumbering{arabic} }
%    \end{macrocode}
% \end{macro}
%
% \begin{macro}{\backmatter}
% 论文后置部分开始命令。
%    \begin{macrocode}
\NewDocumentCommand \backmatter { }
  {
%    \end{macrocode}
% 
% 当双面模式时，确保后置部分内容在奇数页开始。
%    \begin{macrocode}
    \cleardoublepage
%    \end{macrocode}
% 
% 设置章节编号样式。在每一章开始时，节的计数器需要手动设置为 0。
%    \begin{macrocode}
    \setcounter{section}{0}
    \ctexset
      {
        secnumdepth             = 1,
        chapter / numbering     = false,
        chapter / format        += \setcounter{section}{0},
        section / name          = {,.},
        section / number        = {\arabic{section}},
      }
    \bool_if:NF \g_@@_is_bachelor_bool
      { \addtocontents{toc}{\protect\vspace{25bp}} }
  }
%    \end{macrocode}
% \end{macro}
%
% 文档类初始设置。
%    \begin{macrocode}
\keys_set:nn { ecnu / info }
  {
    degree          = academic,
    clc             = {(分类号)},
    secret-level    = {},
    title           = {(中文标题)},
    title-en        = {(English~Title)},
    author          = {(作者姓名)},
    author-en       = {(Author~Name)},
    student-id      = {(学号)},
    department      = {(学院)},
    department-en   = {(Department)},
    major           = {(专业)},
    major-en        = {(Major)},
    supervisor      = {(导师姓名)},
    supervisor-en   = {(Supervisor)},
    academic-title  = {(导师职称)},
    keywords        = {关键词1, 关键词2, 关键词3},
    keywords-en     = {keyword1, keyword2, keyword3}
  }
%    \end{macrocode}
%
%
%    \begin{macrocode}
%</class>
%    \end{macrocode}
%
% \iffalse
%
% \subsubsection{预设文字}
%
%    \begin{macrocode}
%<*definition>
%    \end{macrocode}
% 
% 模板中的一些预设文字，用于封面、声明页、页眉等部分。
%    \begin{macrocode}
\tl_const:Nn \c_@@_info_name_tl               { 姓名 }
\tl_const:Nn \c_@@_info_student_id_tl         { 学号 }
\tl_const:Nn \c_@@_info_student_id_en_tl      { Student ~ ID }
\tl_const:Nn \c_@@_info_department_tl         { 学院 }
\tl_const:Nn \c_@@_info_school_tl             { 院系 }
\tl_const:Nn \c_@@_info_major_tl              { 专业 }
\tl_const:Nn \c_@@_info_major_professional_tl { 专业学位类别 }
\tl_const:Nn \c_@@_info_field_tl              { 研究方向 }
\tl_const:Nn \c_@@_info_field_professional_tl { 专业学位领域 }
\tl_const:Nn \c_@@_info_candidate_tl          { 学位申请人 }
\tl_const:Nn \c_@@_info_supervisor_tl         { 指导教师 }
\tl_const:Nn \c_@@_info_academic_title_tl     { 职称 }
\tl_const:Nn \c_@@_university_code_tl         { 学校代码 }
\tl_const:Nn \c_@@_university_code_en_tl      { University ~ Code }
\tl_const:Nn \c_@@_university_code_val_tl     { 10269 }
\tl_const:Nn \c_@@_year_tl                    { 年 }
\tl_const:Nn \c_@@_month_tl                   { 月 }
\tl_const:Nn \c_@@_day_tl                     { 日 }
\tl_const:Nn \c_@@_date_tl                    { 日期 }
\tl_const:Nn \c_@@_orig_sign_tl               { 承诺人签名 }
\tl_const:Nn \c_@@_author_sign_tl             { 作者签名 }
\tl_const:Nn \c_@@_supervisor_sign_tl         { 导师签名 }
\tl_const:Nn \c_@@_orig_title_ug_tl           { 华东师范大学学位论文诚信承诺 }
\tl_const:Nn \c_@@_auth_title_ug_tl           { 华东师范大学学位论文使用授权说明 }
\tl_const:Nn \c_@@_orig_title_pg_tl           { 华东师范大学学位论文原创性声明 }
\tl_const:Nn \c_@@_auth_title_pg_tl           { 华东师范大学学位论文著作权使用声明 }
\tl_const:Nn \c_@@_orig_text_ug_tl
  {
    \indent 本毕业论文是本人在导师指导下独立完成的，内容真实、可靠。本人
    在撰写毕业论文过程中不存在请人代写、抄袭或者剽窃他人作品、伪造或者
    篡改数据以及其他学位论文作假行为。
    \par
    本人清楚知道学位论文作假行为将会导致行为人受到不授予/撤销学位、开除
    学籍等处理（处分）决定。本人如果被查证在撰写本毕业论文过程中存在学位
    论文作假行为，愿意接受学校依法作出的处理（处分）决定。
  }
\tl_const:Nn \c_@@_auth_text_ug_tl
  {
    \indent 本论文的研究成果归华东师范大学所有，本论文的研究内容不得以其
    它单位的名义发表。本学位论文作者和指导教师完全了解华东师范大学有关
    保留、使用学位论文的规定，即：学校有权保留并向国家有关部门或机构送交
    论文的复印件和电子版，允许论文被查阅和借阅；本人授权华东师范大学可以
    将论文的全部或部分内容编入有关数据库进行检索、交流，可以采用影印、缩印
    或其他复制手段保存论文和汇编本学位论文。
    \par
    保密的毕业论文（设计）在解密后应遵守此规定。
  }
\tl_const:Nn \c_@@_orig_text_pg_part_one_tl
  { 郑重声明：本人呈交的学位论文 }
\tl_const:Nn \c_@@_orig_text_pg_part_two_tl
  {
    是在华东师范大学攻读硕士/博士（请勾选）学位期间，
    在导师的指导下进行的研究工作及取得的研究成果。除文中已经注明引用的内容外，
    本论文不包含其他个人已经发表或撰写过的研究成果。
    对本文的研究做出重要贡献的个人和集体，
    均已在文中作了明确说明并表示谢意。
  }
\tl_const:Nn \c_@@_auth_text_pg_part_one_tl
  {
    系本人在华东师范大学攻读学位期间在导师指导下完成的硕士/博士（请勾选）
    学位论文，本论文的著作权归本人所有。本人同意华东师范大学根据相关规定
    保留和使用此学位论文，并向主管部门和学校指定的相关机构送交学位论文的
    印刷版和电子版；允许学位论文进入华东师范大学图书馆及数据库被查阅、借阅；
    同意学校将学位论文加入全国博士、硕士学位论文共建单位数据库进行检索，
    将学位论文的标题和摘要汇编出版，采用影印、缩印或者其它方式合理复制
    学位论文。
  }
\tl_const:Nn \c_@@_auth_text_pg_part_two_tl
  { 本学位论文属于（请勾选） }
\tl_const:Nn \c_@@_auth_text_pg_option_one_tl
  {
    1. ~ 经华东师范大学相关部门审查核定的“内部”或“涉密”学位论文*，\\
    于 \quad \quad \quad \quad 年 \quad \quad 月 \quad \quad 日
    解密，解密后适用上述授权。
  }
\tl_const:Nn \c_@@_auth_text_pg_option_two_tl
  { 2. ~ 不保密，适用上述授权。 }
\tl_const:Nn \c_@@_auth_others_text_pg_tl
  {
    * “涉密”学位论文应是已经华东师范大学学位管理办公室或
    保密委员会审定过的学位论文（需附获批的《华东师范大学
    研究生申请学位论文“涉密”审批表》方为有效），未经上述
    部门审定的学位论文均为公开学位论文。此声明栏不填写的，
    默认为公开学位论文，均适用上述授权）。
  }
\prop_const_from_keyval:Nn \c_@@_thesis_type_prop
  {
    doctor    = 博士,
    master    = 硕士,
    bachelor  = 学士,
  }
\tl_const:Nn \c_@@_graduation_year_tl         { 届 }
\tl_const:Nn \c_@@_undergraduate_tl           { 本科生 }
\tl_const:Nn \c_@@_postgraduate_tl            { 研究生 }
\tl_const:Nn \c_@@_professional_tl            { 专业 }
\tl_const:Nn \c_@@_degree_tl                  { 学位 }
\tl_const:Nn \c_@@_thesis_tl                  { 论文 }
\tl_const:Nn \c_@@_cic_index_tl               { 分类号 }
\tl_const:Nn \c_@@_secret_level_tl            { 密级 }
\tl_const:Nn \c_@@_thesis_title_tl            { 论文题目 }
\tl_const:Nn \c_@@_university_name_tl         { 华东师范大学 }
\tl_const:Nn \c_@@_university_name_en_tl      { East ~ China ~ Normal ~ University }
\tl_const:Nn \c_@@_professional_en_tl         { Professional }
\tl_const:Nn \c_@@_degree_en_tl               { Degree }
\tl_const:Nn \c_@@_thesis_en_tl               { Thesis }
\tl_const:Nn \c_@@_dissertation_en_tl         { Dissertation }
\tl_const:Nn \c_@@_for_en_tl                  { for }
\tl_const:Nn \c_@@_masters_en_tl              { Master's }
\tl_const:Nn \c_@@_doctoral_en_tl             { Doctoral }
\tl_const:Nn \c_@@_thesis_title_en_tl         { Title }
\tl_const:Nn \c_@@_info_school_en_tl          { Department / School }
\tl_const:Nn \c_@@_info_major_en_tl           { Major }
\tl_const:Nn \c_@@_info_major_pro_en_tl       { Category }
\tl_const:Nn \c_@@_info_field_en_tl           { Research ~ Direction }
\tl_const:Nn \c_@@_info_field_pro_en_tl       { Field }
\tl_const:Nn \c_@@_info_candidate_en_tl       { Candidate }
\tl_const:Nn \c_@@_info_supervisor_en_tl      { Supervisor }
\clist_const:Nn \c_@@_month_en_clist
  {
    January, February, March, April, May, June, July,
    August, September, October, November, December
  }
\tl_const:Nn \c_@@_name_toc_tl                { 目录 }
\tl_const:Nn \c_@@_name_lot_tl                { 表格目录 }
\tl_const:Nn \c_@@_name_lof_tl                { 插图目录 }
\tl_const:Nn \c_@@_name_chinese_abstract_tl   { 中文摘要 }
\tl_const:Nn \c_@@_name_abstract_tl           { 摘要 }
\tl_const:Nn \c_@@_name_abstract_en_tl        { Abstract }
\tl_const:Nn \c_@@_name_keywords_tl           { 关键词 }
\tl_const:Nn \c_@@_name_keywords_en_tl        { Keywords }
\tl_const:Nn \c_@@_colon_tl                   { ： }
\tl_const:Nn \c_@@_colon_en_tl                { : }
\tl_const:Nn \c_@@_left_guillemet_tl          { 《 } 
\tl_const:Nn \c_@@_right_guillemet_tl         { 》 } 
\tl_const:Nn \c_@@_left_round_bracket_tl      { （ } 
\tl_const:Nn \c_@@_right_round_bracket_tl     { ） } 
\tl_const:Nn \c_@@_name_appendix_tl           { 附录 }
\tl_const:Nn \c_@@_name_acknowledgement_tl    { 致谢 }
\tl_const:Nn \c_@@_name_epilogue_tl           { 后记 }
\tl_const:Nn \c_@@_name_pdf_creator_tl        { LaTeX~ with~ ecnuthesis~ class }
%    \end{macrocode}
%
%    \begin{macrocode}
%</definition>
%    \end{macrocode}
%
% \fi
% \end{implementation}